\documentclass[11pt,letterpaper]{article}
\usepackage[margin=0.88in]{geometry}
\usepackage[T1]{fontenc}
\usepackage{lmodern,microtype,amsmath,amssymb,amsthm,mathtools,mathrsfs}
\usepackage{booktabs,longtable,array,xcolor,enumitem,fancyhdr}
\usepackage[hyphens]{url}
\usepackage[colorlinks=true,linkcolor=blue!45!black,citecolor=blue!45!black,urlcolor=blue!45!black]{hyperref}
\usepackage{listings}
\hypersetup{pdftitle={Shared Arithmetic Witnesses for Direct ZFC-Consistency Quartics},pdfauthor={ProveIt research contribution}}
\lstset{basicstyle=\ttfamily\footnotesize,breaklines=true,columns=fullflexible,frame=single}
\setlength{\emergencystretch}{3em}
\setlist{itemsep=3pt,topsep=5pt}
\newcommand{\N}{\mathbb N}
\newcommand{\Z}{\mathbb Z}
\newcommand{\Q}{\mathbb Q}
\newcommand{\Con}{\operatorname{Con}}
\newcommand{\Norm}{\mathcal N}
\newcommand{\Ccal}{\mathcal C}
\newtheorem{theorem}{Theorem}[section]
\newtheorem{lemma}[theorem]{Lemma}
\newtheorem{proposition}[theorem]{Proposition}
\newtheorem{corollary}[theorem]{Corollary}
\theoremstyle{definition}\newtheorem{definition}[theorem]{Definition}
\theoremstyle{remark}\newtheorem{remark}[theorem]{Remark}
\pagestyle{fancy}\fancyhf{}
\fancyhead[L]{\small Shared arithmetic witnesses for ZFC consistency}
\fancyhead[R]{\small ProveIt research contribution}
\fancyfoot[C]{\thepage}
\setlength{\headheight}{14pt}
\allowdisplaybreaks
\title{\textbf{Shared Arithmetic Witnesses\\for Direct ZFC-Consistency Quartics}\\[5pt]
\large An explicit 608-natural-variable equation\\and a 1,685-integer-variable companion}
\author{Prepared for Vladimir Reshetnikov's ProveIt project\\
\small AI-assisted research with independent mathematical and exact algebraic reviews}
\date{23 October 2026}
\begin{document}
\maketitle
\begin{abstract}
We construct an ordinary integer polynomial of exact degree four in 608
nonnegative unknowns whose absence of zeros is provably equivalent in ZFC
to syntactic $\Con(\mathrm{ZFC})$. Its literal specification has 491 quadratic
residuals, 7,582 expanded nonzero monomials, and coefficient height
79,824,571. A degree-four companion in 1,685 unrestricted integer unknowns
has 850 residuals, 12,082 monomials, and height 79,827,802.
The construction uses the existing numerical C222 proof calculus directly;
it does not simulate K215 or instantiate Jones's universal polynomial.

The principal new compression is simultaneous arithmetic certification:
one Pell growth solution, one auxiliary Pell solution, and one Chinese
remainder parameter serve every positive exponentiation in the complete
construction. Equal exponent arguments share their index certificates.
All binomial graphs share one radix, fixed by the sum of their upper
arguments. A self-bounding identity replaces the large power previously
used in each factorial graph. Cofinal beta codes and certified elimination
of intermediate variables complete the reduction.

The 608-variable result reduces the inspected incoming direct construction's
3,614 variables by 83.18\%. It is a new compact point in this direct
quartic branch, not a global minimum in variables or degree. We supply
ordinary unbounded proofs and independent exact checks of the emitted
polynomials. No new end-to-end Lean proof or encoded object-language ZFC
derivation is claimed.
\end{abstract}
\newpage
\tableofcontents
\newpage

\section{Result, meaning, and comparison}

Throughout, $\N=\{0,1,2,\ldots\}$. Consistency means the absence of a
finite syntactic refutation in the fixed ZFC calculus used by ProveIt,
not the existence of a standard model or an additional soundness assumption.
Write its arithmetic proof predicate as $\operatorname{Prf}_{\rm ZFC}(p,\bot)$:
\[
 \Con(\mathrm{ZFC})\;:\!\iff\;
 \neg\exists p\in\N\ \operatorname{Prf}_{\rm ZFC}(p,\bot).
\]

\begin{definition}[The delivered natural polynomial]
Let $r_0,\ldots,r_{490}\in\Z[x_0,\ldots,x_{607}]$ be the literal quadratic
polynomials in \texttt{artifacts/natural/equation.json}, printed completely
in Appendix~\ref{app:literal}. Define
\begin{equation}\label{eq:main}
 \boxed{D_{\N}(x_0,\ldots,x_{607})
      =\sum_{i=0}^{490} r_i(x_0,\ldots,x_{607})^2.}
\end{equation}
The companion \texttt{expanded\_polynomial.json} is the same polynomial
with all products expanded and equal monomials collected.
\end{definition}

\begin{theorem}[Main result]\label{thm:main}
For the fixed polynomial~\eqref{eq:main}, ZFC proves
\[
 \bigl(\exists\mathbf x\in\N^{608}\ D_{\N}(\mathbf x)=0\bigr)
 \;\longleftrightarrow\;
 \exists p\in\N\ \operatorname{Prf}_{\rm ZFC}(p,\bot).
\]
Consequently ZFC proves
\[
 \boxed{\bigl(\forall\mathbf x\in\N^{608}\ D_{\N}(\mathbf x)\ne0\bigr)
         \;\longleftrightarrow\;\Con(\mathrm{ZFC}).}
\]
The unrestricted-integer polynomial $D_{\Z}$ specified in
Section~\ref{sec:integer} satisfies the corresponding theorem over
$\Z^{1685}$. Both polynomials have exact degree four.
\end{theorem}

The proof is assembled in Sections~2--\ref{sec:assembly}. Its classical
arithmetic inputs and inherited proof-calculus calibration are stated
explicitly. The assertion that ZFC proves the equivalence denotes an ordinary
mathematical derivation formalizable in ZFC, including its finite arithmetic
constructions. It does not assert that this archive contains a proof object
in the repository's encoded ZFC calculus. The distinction does not make the
equivalence conditional on a consistency assumption.

\subsection{The measured improvement}

\begin{center}\small
\begin{tabular}{lrrrr}
\toprule
Construction & Unknowns & Residuals & Monomials & Height\\
\midrule
Earlier direct C222, $\N$ &16,036&12,366&700,865&$3.2031\cdot10^{14}$\\
Incoming shared certificates, $\N$ &3,614&2,954&15,996&127,778,545\\
\textbf{This construction, $\N$}&\textbf{608}&\textbf{491}&\textbf{7,582}&\textbf{79,824,571}\\
\textbf{This construction, $\Z$}&\textbf{1,685}&\textbf{850}&\textbf{12,082}&\textbf{79,827,802}\\
\bottomrule
\end{tabular}
\end{center}

Every row has degree four. The first two rows are authenticated comparison
artifacts, not newly proved results of this report
\cite{repo,sharedparent,directreview}. The first height is
320,307,372,009,917 exactly. Relative to the incoming natural parent,
the final equation uses 3,006 fewer variables, 2,463 fewer residuals,
and 8,414 fewer monomials. These are complete-polynomial counts.

The pinned repository already contains nine-positive-variable constructions,
ten-integer-variable constructions, and a 37-integer-variable quartic on
different branches \cite{repo}. This report does not improve their variable
counts. Its contribution is a much smaller direct proof-calculus quartic,
with moderate literal coefficients and a complete finite materialization.
Sun's ten-integer-variable representation theorem is also already represented
in the repository \cite{sun}; it is not a new result here.

There is no canonical scalar definition of simplicity. Variables, degree,
monomials, coefficient height, circuit size, printed description, and proof
dependencies are different costs. We do not claim that 608 is optimal,
that degree four is minimal for this particular consistency predicate, or
that the new sharing lemmas are unprecedented in all literature.
Their concrete use and complete composition are new relative to the
inspected incoming direct construction.

\section{A fixed numerical proof calculus}

\subsection{The inherited logical endpoint}

The source is C222, with truth label $0$ and contradiction label $1$.
The supplied provenance contains the actual source declarations
\begin{lstlisting}
KEndpoint.c222_iff_zf_refutation :
  KCode222.C222 Code254.F <-> BProv ZFax_s [] Form.fBot

ConstructibleChoice.zf_refutation_iff_zfc_refutation :
  BProv ZFax_s [] Form.fBot <-> BProv ZFCax_s [] Form.fBot
\end{lstlisting}
The ASCII arrows stand for the corresponding Lean equivalence symbol.
The first declaration has no machine-execution hypothesis. The second uses
constructible-universe removal of Choice. We inherit these mathematical
calibrations; we neither redefine consistency to mean acceptance by a new
checker nor infer them merely from numerical tests. Their exact source
files and hashes are included in
\texttt{provenance/inherited/CALIBRATION\_PINS.json}.

The audited repository head is
\texttt{8b328dc1801a506cdf448798702192a1d8d96f18}.
The immediate incoming archive retains the calibration pin
\texttt{08200eb10fde2db408c69dffd840d6920a569b1f}.
These pins serve different purposes.

\subsection{Literal arithmetic rules}

The numerical constructors are
\[
 I(a,b)=2(a+b)(a+b+1)+4b,\quad
 A(a,b)=I(a,b)+2,\quad M(a,b)=I(a,b)+1,\quad E(a,b)=I(a,b)+3.
\]
Also put $N(a)=I(a,1)$ and $X(a,b)=N(A(a,N(b)))$.
These constructor symbols are scoped to this section, separately from
the later Pell notation.
The nine axiom-code expressions, in literal order, are
\begin{align*}
 h_1&=X(a,E(a,b)),\\
 h_2&=I(A(a,I(b,c)),I(A(a,b),A(a,c))),\\
 h_3&=I(E(a,b),I(M(c,a),M(c,b))),\\
 h_4&=I(A(a,A(b,c)),A(b,A(a,c))),\\
 h_5&=I(E(c,a),I(M(c,b),M(a,b))),\\
 h_6&=I(X(a,A(a,b)),b),\\
 h_7&=I(E(c,a),I(E(c,b),A(M(a,b),E(a,b)))),\\
 h_8&=I(I(I(c,d),a),I(I(a,c),I(b,c))),\\
 h_9&=\operatorname{Bundle}(a,b,c).
\end{align*}
The last expression has the complete finite definition
\begin{align*}
 \operatorname{Req}(x,k,z)
   &=A(x,I(z,N(A(x,I(z,N(M(x,k))))))),\\
 \operatorname{Sub}(k,x,y)&=A(k,I(M(k,x),M(k,y))),\\
 t_1&=I(\operatorname{Sub}(k,x,y),\operatorname{Req}(y,k,z)),\\
 t_2&=N(A(x,t_1)),\qquad t_3=N(A(x,I(z,t_2))),\\
 \operatorname{Clos}(x,y,k,z)&=A(y,I(M(y,k),t_3)),\\
 b_1&=A(z,I(\operatorname{Clos}(x,y,z,z),N(\operatorname{Req}(y,z,z)))),\\
 b_2&=I(M(x,y),A(z,M(x,y))),\\
 \operatorname{Bundle}(x,y,z)&=N(I(b_2,b_1)).
\end{align*}

Use premise labels $p,q$ and conclusion label $v$. The twelve branch
polynomials are $p^2+q^2+v^2$ (truth), the nine expressions
$p^2+q^2+(v-h_j)^2$, with the additional term
\[
 \left(((c-a)(c-b))^2-(d+1)^2\right)^2
\]
in the ninth, and the two rule expressions
\[
 (q-I(p,v))^2,\qquad q^2+(v-A(a,p))^2.
\]
Let $\mathscr H$ be their product. Over naturals $\mathscr H=0$ says
exactly that one branch holds. The ninth branch's extra equation enforces
$c\ne a,b$, with witness $d=|(c-a)(c-b)|-1$ when the product is nonzero.
All expressions are evaluated using shared additions and multiplications;
there is no unspecified axiom predicate.

\section{Finite proof flow and a cofinal beta code}

\subsection{The bounded carrier}

Fix $t,m\ge1$, $M=(m+1)^t$, and labels $c_i,p_i,q_i\le m$ for $i<t$.
Consider
\[
 P(Z)=\prod_{i<t}(1+c_iZ),\qquad
 Q(Z)=(1+Z)\prod_{i<t}(1+p_iZ)(1+q_iZ).
\]
\begin{lemma}[The inherited bounded-factor carrier]\label{lem:carrier}
$P(M)=Q(M)$ if and only if $P=Q$ as integer polynomials.
\end{lemma}
\begin{proof}
If $\deg Q\ge t+1$, then $Q(M)\ge M^{t+1}$, whereas
\[
 P(M)\le(1+mM)^t<((m+1)M)^t=M^{t+1}.
\]
Equality therefore forces both polynomials to have at most $t$ nonconstant
linear factors. Each has coefficient sum at most $(m+1)^t=M$ and
constant coefficient $1$; every coefficient is strictly below $M$.
If $P-Q$ were nonzero, its least nonzero coefficient would be divisible
by $M$, by dividing the zero evaluation by the appropriate power of $M$
and reducing modulo $M$. Its absolute value is smaller than $M$,
a contradiction.
\end{proof}

Unique factorization in $\Q[Z]$ gives, for each positive label $j$,
\[
 \#\{i:c_i=j\}-\#\{i:p_i=j\}-\#\{i:q_i=j\}=[j=1].
\]
Truth label $0$ contributes a factor $1$. This is a finite nonnegative
flow of valid single-conclusion Horn rules to target $1$.
Such a flow implies derivability: if $1$ were outside the derivable
closure, sum the finitely supported balance outside that closure.
Every rule with its conclusion outside has at least one premise outside,
so each contribution is nonpositive; the target contributes $1$,
a contradiction. Cycles and repeated premises cause no exception.
Conversely a finite proof tree supplies the balances.

\subsection{One integral accumulator}

Encode the flow by
\[
 S_0=M+1,\quad S_t=1,\qquad
 S_{i+1}(1+c_iM)=S_i(1+p_iM)(1+q_iM).
\]
Telescoping and Lemma~\ref{lem:carrier} give the balances. Conversely,
process a proof tree by expanding open obligations; let $S_i$ be the
product of their positive weights $1+jM$. All intermediate values are
natural numbers.

\begin{lemma}[Cofinal beta coding]\label{lem:beta}
Every finite natural trace $S_0,\ldots,S_t$ has $x,y\in\N$ with
\[
 S_i\le y,\qquad x\bmod(1+(i+1)y)=S_i\quad(0\le i\le t).
\]
The scale $y$ can be chosen arbitrarily large.
\end{lemma}
\begin{proof}
Choose $y$ above all trace entries and divisible by $t!$. The moduli
$1+(i+1)y$ are pairwise coprime: a common divisor at indices $i<j$
is coprime to $y$ and divides $j-i$, while $j-i$ divides $y$.
CRT supplies $x$. All prescribed residues are smaller than their moduli.
\end{proof}

A local row therefore needs only the eleven witnesses
\[
 (p,q,v,a,b,c,d,S,S',u,u'),
\]
with label bounds $m$, digit bounds $y$, and common bound
$V=m+y+U$ for the remaining six witnesses. Besides $\mathscr H=0$,
the local equations are
\begin{align*}
 S'(1+vM)&=S(1+pM)(1+qM),\\
 x&=S+(1+(i+1)y)u,\\
 x&=S'+(1+(i+2)y)u'.
\end{align*}
Since $S,S'\le y$, these are exact beta digits. The two deleted
remainder-slack coordinates reconstruct uniquely. Removing them without
the digit bounds would be invalid.

The endpoint equations, with $M=M_0+1$, are
\begin{align*}
 x&=M_0+2+(y+1)u_0,& y&=M_0+r_0+2,\\
 x&=1+(1+(t+1)y)u_t,& 1+(t+1)y&=r_t+2.
\end{align*}
The code chooses $t=t_0+1,m=m_0+1$. A proof of target $1$ can be
chosen with $t,m\ge1$, so these shifts preserve existence.

\section{Eliminating the bounded universal quantifier}

Let the sum of the three local squared residuals and $\mathscr H$ be
the signed polynomial $E-F$, where $E,F$ are the actual nonnegative
arithmetic circuits emitted by \texttt{direct\_calculus.py}.

\begin{lemma}[Heterogeneous CRT certificate]\label{lem:crt}
Let the $j$th local coordinate have bound $u_j\le V$. Suppose
$H>E(i,\mathbf w),F(i,\mathbf w)$ on every allowed local assignment,
$H\ge t$, and $H\ge u_j+1$. Set
\[
 T=H!,\qquad d_i=1+(i+1)T,\qquad \mathcal B=\prod_{i<t}d_i.
\]
Then $\forall i<t\ \exists\mathbf w\le\mathbf u:\ E(i,\mathbf w)=F(i,\mathbf w)$
is equivalent to existence of encoded coordinates $a_j\ge1$ and $C\ge0$
satisfying
\[
 \mathcal B\mid\binom{a_j}{u_j+1},\qquad
 \mathcal B\mid 1+(C+1)T,\qquad
 E(C,\mathbf a)\equiv F(C,\mathbf a)\pmod{\mathcal B}.
\]
\end{lemma}
\begin{proof}
Fix $i<t$ and a prime divisor $\ell$ of $d_i$.
As $\ell\nmid H!$, we have $\ell>H$. The index condition gives
$C\equiv i\pmod\ell$. Multiplying the binomial by $(u_j+1)!$ gives
the falling factorial $\prod_{h=0}^{u_j}(a_j-h)$.
Primality supplies some $w_j\in[0,u_j]$ with
$a_j\equiv w_j\pmod\ell$. Thus $E(i,\mathbf w)\equiv
F(i,\mathbf w)\pmod\ell$; both values are below $\ell$, so they agree.
One decodes modulo a prime divisor, not the generally composite $d_i$.

Conversely choose all finite local witnesses first. The $d_i$ are
pairwise coprime since $T$ is divisible by every positive index
difference. CRT supplies the encoded coordinates and index.
The falling factorial vanishes modulo each $d_i$. Since
$(u_j+1)!\mid T$ and $\gcd(T,\mathcal B)=1$, cancel this factorial
to obtain the stated divisibility. Representatives $a_j$ can be made
positive by adding multiples of $\mathcal B$; the same falling-factorial
argument proves preservation of the binomial condition.
\end{proof}

\subsection{A certified majorant}

For a nonnegative circuit with inputs at most $\Omega\ge4$, assign
exponents by $d(\text{input})=1$, $d(0)=d(1)=0$,
$d(2)=d(3)=d(4)=1$, and
\[
 d(f+g)=1+\max(d(f),d(g)),\qquad d(fg)=d(f)+d(g).
\]
Induction gives $f\le\Omega^{d(f)}$. The supplied 486-node certificate
for the actual $E+F$ circuit gives $D=6,893,164$.
Use
\[
 \Omega=t+x+y+m+M_0+V+4,\qquad H=\Omega^D+t+V+1.
\]
This meets every hypothesis of Lemma~\ref{lem:crt}. The fixed power
$\Omega^D$ is expanded by repeated squaring into multiplication gates.
$D$ is a majorant exponent, not the final polynomial's total degree.
This change may increase witness magnitudes substantially.

\subsection{The product of moduli}

Use the exact arithmetic identity
\begin{align*}
 Q&=(1+(t+1)T)^t+2,\\
 \mathcal B&=\left(T^t\,t!\binom{Q+t}{t}\right)\bmod(TQ-1).
\end{align*}
The numerator is $\prod_{j=1}^{t}(TQ+jT)$, congruent to
$\prod_{j=1}^{t}(1+jT)$ modulo $TQ-1$.
The latter product is positive and smaller than $TQ-1$, because
$T\ge1$ and $Q$ exceeds its elementary factor bound.
Division is encoded by $n=dq+r$, $d=r+s+1$ with natural witnesses.
All remaining variable powers, binomials, and factorials are replaced
by the finite certificates in the next section.

\section{A shared certificate for a finite family of powers}
\label{sec:shared-arithmetic}

The principal change is a change of quantified witnesses.  A compiler that
expands every arithmetic relation independently repeats number-theoretic
infrastructure that can be chosen once for the entire finite system.
We prove simultaneous existence before identifying any witness slots.
The resulting identifications preserve existential solvability, rather
than the set of assignments to the old auxiliary coordinates.
These reductions are new relative to the audited direct-compiler lineage;
we make no claim that every underlying observation is new to the literature.

\subsection{The classical positive-power interface}

For an integer $a\geq2$, let $\chi_a(j),\psi_a(j)$ be the natural Pell
sequences determined by
\begin{align*}
 \chi_a(0)&=1,&\psi_a(0)&=0,\\
 \chi_a(j+1)&=a\chi_a(j)+(a^2-1)\psi_a(j),&
 \psi_a(j+1)&=\chi_a(j)+a\psi_a(j).
\end{align*}
Thus $\chi_a(j)^2-(a^2-1)\psi_a(j)^2=1$.  Write
\[
 \mathcal P(a;x,y)=x^2+y^2-(ay)^2-1.
\]
We use the classical Pell characterization of exponentiation, in the
explicit form proved by \texttt{matiyasevic} and
\texttt{eq\_pow\_of\_pell} in Mathlib~\cite{mathlibpell}.
For $n,k>0$, the graph $m=n^k$ is characterized by natural witnesses
\[
 W,A,T,Z,X,Y,U,V,B,S,R
\]
with $A,B\geq2$, $k\leq Y$, $n,k\leq W$, and
\begin{align}
 \mathcal P(A;X,Y)&=\mathcal P(A;U,V)
 =\mathcal P(B;S,R)=0,\label{eq:sa-three-norms}\\
 V&>0,\qquad Y^2\mid V,\label{eq:sa-v-divides}\\
 B&\equiv1\pmod{4Y},&B&\equiv A\pmod U,\label{eq:sa-b-crt}\\
 S&\equiv X\pmod U,&R&\equiv k\pmod{4Y},\label{eq:sa-index-congruences}\\
 X+Yn&\equiv YA+m\pmod T,\label{eq:sa-power-congruence}\\
 2An&=T+n^2+1,\qquad m<T,\label{eq:sa-t-bound}\\
 \mathcal P(W+1;A,WZ)&=0.\label{eq:sa-majorant-norm}
\end{align}
The first block recovers $X=\chi_A(k)$ and $Y=\psi_A(k)$; the last
block supplies the size bound that makes the power congruence exact.
The displayed congruence moves $Yn$ to the left so that the compiler need
not evaluate a potentially negative intermediate expression.

For clarity, the classical facts needed below are
\begin{gather}
 \psi_a(j)\geq j,\qquad \chi_a(j)\geq a^j,\qquad
 \psi_a(j)\equiv j\pmod{a-1},\label{eq:sa-classical-bounds}\\
 \psi_a(j)^2\mid\psi_a\bigl(j\psi_a(j)\bigr),\qquad
 d\mid e\Longrightarrow\psi_a(d)\mid\psi_a(e).
 \label{eq:sa-classical-divisibility}
\end{gather}
Both sequences are strictly increasing in their index.  Their dependence
on $a$ is polynomial, so
congruent base parameters give congruent values at the same index.
These are the same Pell inputs used by the inherited construction.

\subsection{Orienting the congruence quotients}

A congruence $F\equiv G\pmod H$ is often compiled using two natural
quotients, $F+Hu=G+Hv$.  Canonical completeness witnesses let us orient
the four congruences in the positive-power interface and use one quotient
each.  The orientation is proved for a chosen complete family; it is not
assumed for every old certificate.

\begin{lemma}[Oriented positive-power witnesses]
\label{lem:sa-orientation}
In the completeness direction of the positive-power interface, witnesses
can be chosen satisfying
\begin{align*}
 B&=A+Uq_B,&S&=X+Uq_S,\\
 R&=k+4Yq_R,&X+Yn&=YA+m+Tq_P,
\end{align*}
with $q_B,q_S,q_R,q_P\in\N$.  They also satisfy
$B=1+4Y(q_4+1)$ for a natural $q_4$.
\end{lemma}

\begin{proof}
Choose $X=\chi_A(k)$, $Y=\psi_A(k)$ and an auxiliary even index
$M\geq2$ for which $Y^2\mid\psi_A(M)$.  Put
$U=\chi_A(M)$ and $V=\psi_A(M)$.  The standard coprimality argument
permits a natural CRT witness $B\equiv A\pmod U$,
$B\equiv1\pmod{4Y}$.  Since $U>A$, the residue $A$ lies strictly
between zero and $U$, and therefore $B\geq A$.
This proves the first orientation and $B\geq2$.

Take $S=\chi_B(k)$ and $R=\psi_B(k)$.  Joint induction in the two
recurrences shows that both coordinates are nondecreasing in the base
parameter: their initial values agree and their nonnegative recurrence
coefficients increase with the parameter.  Thus $B\geq A$ gives
$S\geq X$.  Also $R\geq k$.  The two corresponding congruences now
have nonnegative difference quotients.

For the remaining orientation, the recurrence gives
\[
 \chi_A(k)-A\psi_A(k)=-\psi_A(k-1)\qquad(k\geq1).
\]
Consequently
\[
 L_0:=X+nY-AY=n\psi_A(k)-\psi_A(k-1)>0.
\]
The power congruence says $L_0\equiv m\pmod T$, while the canonical
choice has $0<m<T$.  If $L_0=m+Tq$ with an integer $q\leq-1$, then
$L_0\leq m-T<0$, a contradiction.  Hence $q=q_P\geq0$.
Finally, $(B-1)/(4Y)$ is a positive integer, since a zero quotient
would give $B=1$; subtracting one defines $q_4\in\N$.
\end{proof}

Each oriented equality immediately implies its old congruence, so
soundness is preserved.  In the old two-quotient convention the lift is
simply $(u,v)=(0,q)$.  Two further bounds need no independent slack.
First, $Y\geq k>0$ and $B=1+4Y(q_4+1)$ imply $B\geq5$.
Second, the equations $T=m+g+1$ and $2An=T+n^2+1$ imply $A\geq2$:
if $A\leq1$, then
\[
 2An\leq2n\leq n^2+1<T+n^2+1,
\]
contradicting the defining equality.  This inference precedes the use
of any Pell theorem and does not assume that the candidate output is
already the correct power.

\subsection{Simultaneous choice of the majorant and auxiliary pair}

\begin{theorem}[One Pell infrastructure for a finite family]
\label{thm:sa-shared-pell}
Let $(n_j,k_j)$, $1\leq j\leq r$, be a nonempty finite family with
$n_j,k_j>0$.  The conjunction $m_j=n_j^{k_j}$ admits the preceding
positive-power certificates with a common $W,A,Z,U,V,B$ and a common
$q_B$.  One may impose the exact equation
\[
 W=\sum_{j=1}^r(n_j+k_j).
\]
Calls with the same exponent may additionally share their complete
Pell-index block.
\end{theorem}

\begin{proof}
We prove completeness by an explicit order of choices.  Fix the intended
outputs and put $W=\sum_j(n_j+k_j)$.  Then $W\geq n_j,k_j$ and
$W\geq2$.  Choose
\[
 A=\chi_{W+1}(W),\qquad WZ=\psi_{W+1}(W).
\]
The ordinate is divisible by $W$ by
\eqref{eq:sa-classical-bounds}, so $Z$ is natural.  The common last
norm holds, and for every $j$,
\[
 n_j^{k_j}\leq n_j^W<(W+1)^W\leq A.
\]
The classical power-bound lemma therefore gives
\[
 0<m_j<T_j:=2An_j-n_j^2-1.
\]
No property of the original choice $\max(n_j,k_j)$, beyond its two
upper bounds, was needed.

Index the distinct supplied exponents by $e_1,\ldots,e_s$ and set
\[
 X_\ell=\chi_A(e_\ell),\quad Y_\ell=\psi_A(e_\ell),\quad
 M=2\prod_{\ell=1}^s e_\ell Y_\ell,
 \quad U=\chi_A(M),\quad V=\psi_A(M).
\]
Every $e_\ell Y_\ell$ divides $M$, so
\eqref{eq:sa-classical-divisibility} gives $Y_\ell^2\mid V$ for
every $\ell$.  Also $M\geq2$ and $V>0$.
The recurrence $\psi_a(j+2)=2a\psi_a(j+1)-\psi_a(j)$ shows
$\psi_a(j)\equiv j\pmod2$, and hence $V$ is even.
The Pell norm implies $\gcd(U,V)=1$.  Thus $U$ is odd and coprime
to every $Y_\ell$, whence
\[
 \gcd\left(U,4\prod_{\ell=1}^sY_\ell\right)=1.
\]
Since $M\geq2$, strict index growth gives $U>A$.

Choose one nonnegative CRT witness
\[
 B\equiv A\pmod U,
 \qquad B\equiv1\pmod{4\prod_\ell Y_\ell}.
\]
It satisfies $B\geq A$ and all the required individual congruences
$B\equiv1\pmod{4Y_\ell}$.  Define
$S_\ell=\chi_B(e_\ell)$ and $R_\ell=\psi_B(e_\ell)$.
Lemma~\ref{lem:sa-orientation} supplies all the natural oriented
quotients.  In particular, each positive quotient $V/Y_\ell^2$ gives
the row $V=Y_\ell^2(v_\ell+1)$ with $v_\ell\in\N$.
The same index block can be supplied to every call having exponent
$e_\ell$.  The output gap and the power-congruence quotient remain
specific to its base and output.

For soundness, take any natural solution of the shared system.
The equation for $W$ gives all the required upper bounds.  The positive
output gap and the $T_j$ equation establish $A\geq2$ as above.
For each individual call, reuse the common $W,A,Z,U,V,B$ in the
classical interface.  Its own index block proves $Y\geq k>0$,
$V>0$, $Y^2\mid V$ and $B\geq2$, and its oriented rows imply the
old congruences.  The classical theorem gives $m_j=n_j^{k_j}$.
Soundness does not require an arbitrary shared auxiliary pair to arise
from the particular product index $M$ used in the completeness proof.
\end{proof}

The product defining $M$ and the CRT choice of $B$ occur only in this
proof.  They are not variable-exponent expressions or additional
operations in the final polynomial.  The final system stores the
common norm and congruence rows and the individual finite certificates.
The released implementation shares exponent blocks only for literally
identical supplied atoms; it has twenty such blocks.  It does not merge
different expressions merely because they might have equal values on
some solutions.

\subsection{Quadratic rows and positive-input specialization}

The norm $x^2-(a^2-1)y^2=1$ can be compiled with a single product
register $h=ay$ and the quadratic residual
\[
 x^2+y^2-h^2-1=0.
\]
The witness $h$ is uniquely determined and natural.  In the main norm
the same $AY$ product is already needed by the power-congruence row.
Squares that occur as terms of a quadratic residual need no separate
register.  Likewise $Y=k+y_0$ and $T=m+g+1$ are exact affine
parametrizations of the two bounds, with explicit inverse maps
$y_0=Y-k$ and $g=T-m-1$ on old solutions.

These changes apply to the positive branch.  A general exponentiation
graph must still handle $k=0$, where $m=1$, and $n=0<k$, where $m=0$.
The present source establishes positive bases and exponents at every
emitted call, so those branches are unnecessary here.  This is a property
of the full construction, not an assumption that arbitrary power inputs
are positive.  In particular the nested calls introduced below remain
positive before their candidate factorial outputs have been identified.
An unguarded norm is quadratic in the form just displayed; multiplying
it by a Boolean selector would instead be cubic, so the simplification
is not asserted for an unchanged guarded implementation.

\section{One radix for all binomial coefficients}
\label{sec:shared-radix}

Suppose a finite family of binomial graphs has natural upper arguments
$N_i$ and lower arguments $K_i$.  The independent compiler used a
separate radix $2^{N_i}+1$ for each call.  Set instead
\[
 L=\sum_i N_i,\qquad b=2^L,\qquad q=b+1.
\]
Here $b=2^L$ is one power graph, compiled by the preceding quadratic
certificate; it is not retained as an exponentiation operator.

\begin{lemma}[Common-radix digit extraction]
\label{lem:sa-common-radix}
For every $i$, the coefficient $c_i=\binom{N_i}{K_i}$ is characterized
by natural $h_i,r_i,u_i,v_i$ satisfying
\begin{align*}
 (q+1)^{N_i}&=q^{K_i}(qh_i+c_i)+r_i,\\
 q^{K_i}&=r_i+u_i+1,\\
 q&=c_i+v_i+1.
\end{align*}
The powers in these displayed semantic equations can all be replaced
by the finite positive-power graphs when their inputs are positive.
\end{lemma}

\begin{proof}
Since $L\geq N_i$, we have $q>2^{N_i}$.  The binomial coefficients
are nonnegative and their sum is $2^{N_i}$, so every coefficient in
\[
 (q+1)^{N_i}=\sum_{j=0}^{N_i}\binom{N_i}{j}q^j
\]
is smaller than $q$.  This is therefore its base-$q$ expansion without
carries.  Its digit in position $K_i$ is $\binom{N_i}{K_i}$, with the
usual value zero for $K_i>N_i$.
The first two rows perform Euclidean division by $q^{K_i}$; the last
row restricts the next remainder to the interval $[0,q)$.
Together they select exactly that digit.  Conversely Euclidean
division supplies $h_i,r_i$ and both nonnegative slacks.
\end{proof}

The digit lemma itself also handles $N_i=0$ or $K_i=0$ under the
convention $q^0=1$.  In the positive-only released construction the
upper and lower arguments are positive, so $L>0$ and all power calls
remain in their supported branch.  Equal lower arguments now use the
same denominator power $q^{K_i}$; equal upper arguments use the same
numerator power.  Literal cache equality, rather than an unproved
algebraic simplification, determines these additional identifications.

The distinction between $L$ and the Pell majorant $W$ is essential.
The latter bounds the inputs of all power calls, including bases $q$
and $q+1$.  Setting $q=2^W+1$ would create a circular size requirement.
Our order of choices is different: fix the intended high-level arithmetic
values, then their binomial upper arguments and $L$; choose $q=2^L+1$;
then all power bases and exponents are fixed, permitting the subsequent
choice of $W,A,Z,U,V,B$.

For soundness, the shared Pell system first establishes every power
graph pointwise.  The literal equation $L=\sum_iN_i$ then gives the
bound for every digit extraction.  Thus the binomial graphs are proved
before using any factorial correctness claim.  In completeness, even
when a factorial output appears in an upper argument, we first assign
its intended value and then construct these arithmetic certificates.
No upper argument refers to a private Pell witness.

\section{A self-bounding factorial certificate}
\label{sec:self-bounding-factorial}

The inherited factorial macro chose a radius using a further large
power.  Its only purpose was to make a rational approximation close
enough to the integer $n!$.  The candidate factorial output itself
provides a sufficient radius, provided both implications are proved.

\begin{theorem}[Self-bounding factorial]
\label{thm:sa-factorial}
For $n,m\in\N$, put
\[
 Q=mn^2,\qquad C=\binom{Q+n}{n},\qquad P=(Q+1)^n.
\]
Then
\[
 m=n!\quad\Longleftrightarrow\quad
 m=\left\lceil\frac{P}{C}\right\rceil.
\]
Equivalently, $m=n!$ holds precisely when natural $r,s$ satisfy
\begin{equation}
 P+C=mC+r+1,\qquad C=r+s+1.
 \label{eq:sa-factorial-rows}
\end{equation}
\end{theorem}

\begin{proof}
Always $P,C\geq1$.  The two rows are equivalent to
\[
 P=C(m-1)+(r+1),\qquad 1\leq r+1\leq C,
\]
and hence to $m-1<P/C\leq m$.  They force $m\geq1$: if $m=0$,
the first row would give $P+C=r+1\leq C$, contradicting $P\geq1$.
Conversely the ceiling inequalities give the natural witnesses
$r=P+C-mC-1$ and $s=C-r-1$.

When $n=0$, one has $Q=0$ and $P=C=1$, so both sides assert $m=1$.
Now assume $n\geq1$ and $m\geq1$, so $Q=mn^2>0$.  Write $F=n!$
and use the exact rational identity
\[
 R:=\frac{P}{C}=F\prod_{i=1}^{n}\frac{Q+1}{Q+i}.
\]
Every factor lies in $(0,1]$, giving $0<R\leq F$.
Let $d_i=(i-1)/(Q+i)$.  These numbers lie in $[0,1)$, and
\[
 \sum_{i=1}^n d_i
 \leq\frac{n(n-1)}{2Q}
 <\frac{n^2}{2Q}=\frac{1}{2m}.
\]
Induction gives the elementary inequality
$\prod_i(1-d_i)\geq1-\sum_i d_i$.  Therefore
\begin{equation}
 R>F\left(1-\frac{1}{2m}\right).
 \label{eq:sa-factorial-ratio}
\end{equation}

For soundness, suppose $m=\lceil R\rceil$.  Since $F$ is an integer
and $R\leq F$, we have $m\leq F$.  On the other hand $R\leq m$
and \eqref{eq:sa-factorial-ratio} imply
\[
 F<\frac{m}{1-1/(2m)}
   =m+\frac{m}{2m-1}\leq m+1.
\]
The first inequality is strict, including when $m=1$.  Integrality
therefore gives $F\leq m$, and so $F=m$.

For completeness, set $m=F$.  The same estimate gives
$R>F-\tfrac12>F-1$, while $R\leq F$.  Thus $\lceil R\rceil=F$,
and the displayed natural remainder witnesses exist.
\end{proof}

The proof addresses the apparent circularity directly: $m$ is initially
an arbitrary candidate; the fixed-point condition forces it to be the
factorial.  Merely replacing the radius by a smaller expression is not
safe.  For example, the choice $Q=mn$ has a false fixed point at
$n=4,m=23$, although $4!=24$.  The quadratic radius and its strict
estimate are doing mathematical work.

In the compiler, allocate the output $m$ and replace the old radius
power by the polynomial expression $mn^2$.  Keep the remaining
binomial graph, numerator power and rows
\eqref{eq:sa-factorial-rows}.  For a positive factorial input $n$,
one has $Q+n\geq n>0$ and $Q+1\geq1$ even before proving $m=n!$.
Hence the nested binomial upper argument, lower argument, numerator
base and numerator exponent are positive without using the desired
factorial conclusion.  The power and binomial graph proofs may therefore
be applied first, followed by Theorem~\ref{thm:sa-factorial}.

Finally, all binomial coefficients, ceilings, Pell sequences, divisibility
symbols and powers with variable exponents in these sections belong to
the correctness proof and the semantic description of macros.  The
released equation contains only a fixed finite list of integer quadratic
residuals.  Their sum of squares is an ordinary polynomial of degree
four with fixed integer coefficients; no such semantic operation remains
in that polynomial.


\section{Quadratic compilation and safe elimination}\label{sec:compile}

Each arithmetic operation receives its monic gate equation
$g-a-b=0$ or $g-ab=0$. An unguarded Pell norm
$x^2-(a^2-1)y^2=1$ is compiled economically:
\[
 h=ay,\qquad x^2+y^2-h^2-1=0.
\]
Only one product coordinate is needed. This relies on unguarded
residuals: multiplying the last row by a selector would make it cubic.
All power inputs in this construction are independently proved positive.

\begin{lemma}[Natural affine elimination]\label{lem:affine}
If $g=L$ is a monic gate and $L$ is affine with nonnegative integer
coefficients in earlier coordinates, substituting $L$ for $g$ and
removing its defining row preserves quadratic degree and gives a
bijection of natural zero sets.
\end{lemma}
\begin{proof}
Projection has the unique inverse $g=L\ge0$. Affine substitution
preserves degree at most two. Iterate in topological order.
The implementation's support cap of four is a sparse-size choice only.
\end{proof}

\begin{lemma}[Terminal nonlinear elimination]\label{lem:terminal}
Suppose a residual is $g-Q$, where $g$ is absent from $Q$, $Q$ has
nonnegative coefficients and degree at most two, and every occurrence
of $g$ in every other residual is linear. Substituting $Q$ for $g$ and
deleting $g-Q$ preserves quadratic degree and gives a bijection of
natural zero sets.
\end{lemma}
\begin{proof}
Every other row is $a g+P$, with $g$ absent from $P$, and becomes
$aQ+P$, still quadratic. Restore $g=Q\ge0$ uniquely.
Reverse the elimination order to restore a sequence of such coordinates;
each recipe refers only to coordinates present at its deletion stage.
\end{proof}

The raw system has 1,541 variables and 1,424 residuals, including
1,182 triangular nonnegative gate definitions. Removing 724 affine
coordinates leaves 817 variables and 700 residuals. A further 209
terminal eliminations leave 608 variables and 491 residuals.
Every retained coordinate is used. The independent checker derives
each recipe from its defining row, checks all consumers, and does not
import the selection heuristic.

Squaring and summing gives $D_{\N}$, with 7,582 collected monomials.
A final residual has a nonzero quadratic homogeneous part.
The sum of squares of all such parts cannot vanish identically over
$\mathbb R$, so the degree is exactly four.

\section{An unrestricted-integer quartic}\label{sec:integer}

Among the 608 variables, 359 are original nongate inputs and 249
are retained arithmetic gates. Let $I$ be the exact set of 359 input
indices printed in Appendix~\ref{app:integer}. Introduce three
integer variables $a_j,b_j,c_j$ for every $j\in I$, and set
\begin{equation}\label{eq:integer}
 \boxed{D_{\Z}=D_{\N}+
  \sum_{j\in I}(8x_j+3-a_j^2-b_j^2-c_j^2)^2.}
\end{equation}
This is an ordinary polynomial in $608+3\cdot359=1,685$ integer variables.

\begin{proposition}
$D_{\Z}$ has an integer zero if and only if $D_{\N}$ has a natural zero.
\end{proposition}
\begin{proof}
For a natural zero, Legendre's three-square theorem supplies the guards:
$8x_j+3$ cannot be $4^a(8b+7)$.
Conversely, the guards force the 359 original inputs to be nonnegative.
First reverse both elimination stages as polynomial substitutions over
$\Z$, without assuming the retained gates are natural. This restores
the raw integer system. Induction through all 1,182 positive triangular
gate definitions makes every raw gate nonnegative. Thus the restored
assignment is a natural solution, whose projection is a natural zero
of $D_{\N}$.
\end{proof}
The order matters: nonnegativity of the 249 unguarded gates is a
conclusion of their restored equations. The integer polynomial has
850 residuals, 12,082 monomials, and height 79,827,802.

\section{Assembly inside ZFC and exact verification}\label{sec:assembly}

Starting with a natural zero, reverse the eliminations to obtain the
raw arithmetic equations. The shared Pell theorem identifies every
power output. Positivity of all inputs is checked in actual macro order.
The common-radix lemma identifies the binomials; the self-bounding
factorial theorem then identifies the factorial outputs. No positivity
of an uncalibrated factorial output is assumed.
Division and the product identity recover the CRT certificate.
Lemma~\ref{lem:crt} supplies local rule instances, whose exact beta
digits form one accumulator trace. The carrier and flow arguments
yield C222 derivability of $1$. The inherited logical endpoint gives
a syntactic ZFC refutation.

Conversely a finite refutation gives a finite C222 proof and an
accumulator trace. Choose its carrier, then its cofinal beta code,
then a bound on all remaining local witnesses. Form the factorial
CRT moduli and encoded columns. Assign all semantic factorial and
binomial values, then choose the common digit radix from the sum of
the binomial upper arguments. This fixes all power inputs.
Choose the shared Pell witnesses last. The resulting raw assignment
projects to a zero of $D_{\N}$; the integer companion follows.

These arguments use finite syntax, finite induction, integer
polynomial arithmetic, finite CRT, and the stated classical
number-theory theorems. ZFC proves them about its internal natural
numbers. There is no restriction to externally standard codes.
This proves Theorem~\ref{thm:main}; negating its solvability equivalence
gives the requested consistency equivalence. It proves neither
consistency nor inconsistency.

\subsection{Independent checks}

Separate constructor and checker implementations cover:
\begin{itemize}
\item all 25 power graphs, 20 shared index certificates,
      14 binomial graphs, two factorial graphs, and 221 literal
      semantic interface rows;
\item the actual separate global sums for the Pell and digit bounds,
      and dependency checks excluding circular dependence on private
      Pell witnesses;
\item the changed numerical source, all 486 majorant nodes, and the
      complete raw 1,424-row artifact;
\item all 724 affine and 209 terminal reconstructions, including
      218 exact consumer substitutions;
\item every integer guard and every expanded coefficient in both
      final polynomials.
\end{itemize}
Normal Python and \texttt{python -O} executions pass. Eight corruption
controls reject altered coefficients, nonunit definitions, nonlinear
consumers, incorrect restorations, missing guards, and wrong guard
constants. Independent fixtures also test the new arithmetic lemmas.
These checks validate implementations and finite identities; the
unbounded correctness arguments are the proofs and reviews.

\subsection{Reproduction and artifact conventions}

Python 3 and its standard library suffice. From the archive root:
\begin{lstlisting}
python3 verify_all.py
python3 -O verify_all.py
python3 src/build.py --output /tmp/con-zfc-natural
python3 src/integer_export.py --input /tmp/con-zfc-natural
pdflatex -interaction=nonstopmode -halt-on-error report.tex
pdflatex -interaction=nonstopmode -halt-on-error report.tex
\end{lstlisting}
The verifier regenerates the defining outputs in a temporary directory
and compares them byte for byte. A monomial record
\texttt{[c,[i,j,j]]} denotes $c x_i x_j^2$; an empty index list denotes
the constant monomial. Indices are zero-based.
Abbreviations in the proofs are not additional variables.

A SHA-256 manifest identifies the distributed files. The theorem is
tied to the printed residual list and machine-readable literals,
not merely a changing script. No remote repository has been modified
or published by this contribution.

\subsection{Limits and the next question}

The inspected direct parent was already a fully materialized ordinary
construction. This work improves its complete dimensions.
The remaining limitation is dimensional: 608 naturals is far above
the few-variable branches. Shared Pell indices and CRT parameters
may be enormous; the source majorant can increase witness sizes
further. Small equation size is not a practical search method.

The reusable result is that a finite family of powers can share
its large-parameter Pell infrastructure. A further target is a
more variable-efficient interface for the twenty distinct indices,
or a cheaper direct encoding of the eleven local flow coordinates.
Such a change must retain a simultaneous-existence proof. It is
not supplied or counted here.

\begingroup\small\raggedright
\begin{thebibliography}{9}
\bibitem{repo}
V. Reshetnikov and ProveIt contributors.
\emph{DiophantineConsistency: size research and status}.
Snapshot \texttt{8b328dc1801a506cdf448798702192a1d8d96f18}.
\url{https://github.com/VladimirReshetnikov/ProveIt}.
Frozen comparison files are in \texttt{provenance/}.
\bibitem{sharedparent}
ProveIt incoming contribution.
\emph{Shared certificates},
\texttt{ConZFC\_Shared\_Certificates\_2026-10-24.zip},
in the pinned snapshot's \texttt{docs/incoming}.
Its SHA-256 is in \texttt{provenance/SOURCE\_PINS.json}.
The original filename's date is preserved.
\bibitem{directreview}
ProveIt contributors.
\emph{Ordinary correctness review: direct C222 quartic, round 12}.
\texttt{PROPOSAL\_REVIEW\_DIRECT\_C222\_ROUND12.md}
in the pinned DiophantineConsistency directory.
\bibitem{calibration}
ProveIt contributors.
\emph{Code254, KCode222, KEndpoint, and ConstructibleChoice}.
Source pin \texttt{08200eb10fde2db408c69dffd840d6920a569b1f}.
Exact interfaces and hashes are in \texttt{provenance/inherited/}.
\bibitem{mathlibpell}
M. Carneiro and Mathlib contributors.
\emph{Pell's equation and Matiyasevic's theorem},
\texttt{Mathlib/NumberTheory/PellMatiyasevic.lean}.
Especially \texttt{matiyasevic}, \texttt{eq\_pow\_of\_pell},
and the Pell divisibility and congruence lemmas.
Inspected blob \texttt{db9c757dd7cac29260397150fd190d72715ea3f8}.
\url{https://leanprover-community.github.io/mathlib4_docs/Mathlib/NumberTheory/PellMatiyasevic.html}.
The source and Apache license are included.
\bibitem{sun}
Z.-W. Sun.
\emph{Ten unknowns for Hilbert's tenth problem over the integers}.
arXiv:2609.01594v2, 2 September 2026.
\url{https://arxiv.org/abs/2609.01594}.
Used for comparison and its statement of the classical Gauss--Legendre
theorem; not as this construction's compiler.
\end{thebibliography}
\endgroup

\appendix
\section{Complete natural residual list}\label{app:literal}

This is the literal definition of every $r_i$ in
$D_{\N}=\sum_{i=0}^{490}r_i^2$. All $x_0,\ldots,x_{607}$ range over
$\N$. A continuation line continues the same polynomial.
\begingroup
\footnotesize
\setlength{\jot}{2pt}
\begin{align*}
r_{0} &= -x_{5} x_{8} - x_{8} - x_{8} x_{17} + x_{18} \\
r_{1} &= 2\,x_{5} + x_{5}^{2} + 2\,x_{5} x_{17} + x_{14}^{2} \\
&\quad{}+ 2\,x_{17} + x_{17}^{2} - x_{18}^{2} \\
r_{2} &= -x_{12} x_{16} + x_{19} \\
r_{3} &= -1 + x_{15}^{2} + x_{16}^{2} - x_{19}^{2} \\
r_{4} &= -5 - 4\,x_{5} - 4\,x_{5} x_{20} + x_{12} \\
&\quad{}- 4\,x_{17} - 4\,x_{17} x_{20} - 4\,x_{20} \\
r_{5} &= -1 - 2\,x_{5} - x_{5}^{2} - 2\,x_{5} x_{17} \\
&\quad{}- 2\,x_{17} - x_{17}^{2} + x_{22} \\
r_{6} &= x_{11} - x_{21} x_{22} - x_{22} \\
r_{7} &= -x_{10} x_{23} - x_{14} + x_{15} \\
r_{8} &= -1 - x_{5} - 4\,x_{5} x_{24} + x_{16} \\
&\quad{}- 4\,x_{17} x_{24} - 4\,x_{24} \\
r_{9} &= 2 + 2\,x_{5} - x_{6} - x_{6} x_{26} \\
&\quad{}+ x_{14} + 2\,x_{17} - x_{18} - x_{25} x_{26} \\
&\quad{}- x_{26} \\
r_{10} &= -6 - x_{6} + 4\,x_{8} - x_{25} \\
r_{11} &= -x_{0} x_{8} - x_{8} - x_{8} x_{31} + x_{32} \\
r_{12} &= 2\,x_{0} + x_{0}^{2} + 2\,x_{0} x_{31} + x_{28}^{2} \\
&\quad{}+ 2\,x_{31} + x_{31}^{2} - x_{32}^{2} \\
r_{13} &= -x_{12} x_{30} + x_{33} \\
r_{14} &= -1 + x_{29}^{2} + x_{30}^{2} - x_{33}^{2} \\
r_{15} &= -5 - 4\,x_{0} - 4\,x_{0} x_{34} + x_{12} \\
&\quad{}- 4\,x_{31} - 4\,x_{31} x_{34} - 4\,x_{34} \\
r_{16} &= -1 - 2\,x_{0} - x_{0}^{2} - 2\,x_{0} x_{31} \\
&\quad{}- 2\,x_{31} - x_{31}^{2} + x_{36} \\
r_{17} &= x_{11} - x_{35} x_{36} - x_{36} \\
r_{18} &= -x_{10} x_{37} - x_{28} + x_{29} \\
r_{19} &= -1 - x_{0} - 4\,x_{0} x_{38} + x_{30} \\
&\quad{}- 4\,x_{31} x_{38} - 4\,x_{38} \\
r_{20} &= 2 + 2\,x_{0} + x_{0} x_{3} + x_{3} \\
&\quad{}+ x_{3} x_{31} - x_{27} - x_{27} x_{40} + x_{28} \\
&\quad{}+ 2\,x_{31} - x_{32} - x_{39} x_{40} - x_{40} \\
r_{21} &= -6 - 4\,x_{3} - x_{3}^{2} + 2\,x_{3} x_{8} \\
&\quad{}+ 4\,x_{8} - x_{27} - x_{39} \\
r_{22} &= -1 - x_{4} + x_{27} \\
r_{23} &= -x_{0} x_{2} - 2\,x_{2} + x_{45} \\
r_{24} &= -2 + x_{1} - x_{2} x_{41} - x_{4} \\
&\quad{}- x_{41} \\
r_{25} &= -2 + x_{2} - x_{4} - x_{42} \\
r_{26} &= -1 + x_{1} - x_{43} - x_{43} x_{45} \\
r_{27} &= -1 - x_{44} + x_{45} \\
r_{28} &= -3 - x_{0} - x_{1} - 2\,x_{2} \\
&\quad{}- 2\,x_{3} - x_{4} - x_{46} + x_{59} \\
r_{29} &= -16 - 8\,x_{59} - x_{59}^{2} + x_{60} \\
r_{30} &= -x_{60}^{2} + x_{61} \\
r_{31} &= -x_{61}^{2} + x_{62} \\
r_{32} &= -x_{61} x_{62} + x_{63} \\
r_{33} &= -x_{62}^{2} + x_{64} \\
r_{34} &= -x_{64}^{2} + x_{65} \\
r_{35} &= -x_{63} x_{65} + x_{66} \\
r_{36} &= -x_{65}^{2} + x_{67} \\
r_{37} &= -x_{66} x_{67} + x_{68} \\
r_{38} &= -x_{67}^{2} + x_{69} \\
r_{39} &= -x_{69}^{2} + x_{70} \\
r_{40} &= -x_{70}^{2} + x_{71} \\
r_{41} &= -x_{68} x_{71} + x_{72} \\
r_{42} &= -x_{71}^{2} + x_{73} \\
r_{43} &= -x_{72} x_{73} + x_{74} \\
r_{44} &= -x_{73}^{2} + x_{75} \\
r_{45} &= -x_{74} x_{75} + x_{76} \\
r_{46} &= -x_{75}^{2} + x_{77} \\
r_{47} &= -x_{77}^{2} + x_{78} \\
r_{48} &= -x_{76} x_{78} + x_{79} \\
r_{49} &= -x_{78}^{2} + x_{80} \\
r_{50} &= -x_{80}^{2} + x_{81} \\
r_{51} &= -x_{81}^{2} + x_{82} \\
r_{52} &= -x_{79} x_{82} + x_{83} \\
r_{53} &= -x_{82}^{2} + x_{84} \\
r_{54} &= -x_{84}^{2} + x_{85} \\
r_{55} &= -x_{85}^{2} + x_{86} \\
r_{56} &= -x_{83} x_{86} + x_{87} \\
r_{57} &= -x_{86}^{2} + x_{88} \\
r_{58} &= -x_{88}^{2} + x_{89} \\
r_{59} &= -x_{87} x_{89} + x_{90} \\
r_{60} &= -x_{89}^{2} + x_{91} \\
r_{61} &= -2 - x_{0} - x_{2} - x_{3} \\
&\quad{}- x_{46} - x_{90} x_{91} + x_{92} \\
r_{62} &= -1 - 2\,x_{92} - x_{92}^{2} + x_{94} \\
r_{63} &= -x_{93} x_{94} + x_{95} \\
r_{64} &= -x_{8} - x_{8} x_{92} - x_{8} x_{95} - x_{8} x_{101} \\
&\quad{}+ x_{102} \\
r_{65} &= 2\,x_{92} + x_{92}^{2} + 2\,x_{92} x_{95} + 2\,x_{92} x_{101} \\
&\quad{}+ 2\,x_{95} + x_{95}^{2} + 2\,x_{95} x_{101} + x_{98}^{2} \\
&\quad{}+ 2\,x_{101} + x_{101}^{2} - x_{102}^{2} \\
r_{66} &= -x_{12} x_{100} + x_{103} \\
r_{67} &= -1 + x_{99}^{2} + x_{100}^{2} - x_{103}^{2} \\
r_{68} &= -5 + x_{12} - 4\,x_{92} - 4\,x_{92} x_{104} \\
&\quad{}- 4\,x_{95} - 4\,x_{95} x_{104} - 4\,x_{101} - 4\,x_{101} x_{104} \\
&\quad{}- 4\,x_{104} \\
r_{69} &= -1 - 2\,x_{92} - x_{92}^{2} - 2\,x_{92} x_{95} \\
&\quad{}- 2\,x_{92} x_{101} - 2\,x_{95} - x_{95}^{2} - 2\,x_{95} x_{101} \\
&\quad{}- 2\,x_{101} - x_{101}^{2} + x_{106} \\
r_{70} &= x_{11} - x_{105} x_{106} - x_{106} \\
r_{71} &= -x_{10} x_{107} - x_{98} + x_{99} \\
r_{72} &= -1 - x_{92} - 4\,x_{92} x_{108} - x_{95} \\
&\quad{}- 4\,x_{95} x_{108} + x_{100} - 4\,x_{101} x_{108} - 4\,x_{108} \\
r_{73} &= 2 + x_{6} + x_{6} x_{92} + x_{6} x_{95} \\
&\quad{}+ x_{6} x_{101} + 2\,x_{92} + 2\,x_{95} - x_{97} \\
&\quad{}- x_{97} x_{110} + x_{98} + 2\,x_{101} - x_{102} \\
&\quad{}- x_{109} x_{110} - x_{110} \\
r_{74} &= -6 - 4\,x_{6} - x_{6}^{2} + 2\,x_{6} x_{8} \\
&\quad{}+ 4\,x_{8} - x_{97} - x_{109} \\
r_{75} &= -x_{8} - x_{8} x_{92} - x_{8} x_{115} + x_{116} \\
r_{76} &= 2\,x_{92} + x_{92}^{2} + 2\,x_{92} x_{115} + x_{112}^{2} \\
&\quad{}+ 2\,x_{115} + x_{115}^{2} - x_{116}^{2} \\
r_{77} &= -x_{12} x_{114} + x_{117} \\
r_{78} &= -1 + x_{113}^{2} + x_{114}^{2} - x_{117}^{2} \\
r_{79} &= -5 + x_{12} - 4\,x_{92} - 4\,x_{92} x_{118} \\
&\quad{}- 4\,x_{115} - 4\,x_{115} x_{118} - 4\,x_{118} \\
r_{80} &= -1 - 2\,x_{92} - x_{92}^{2} - 2\,x_{92} x_{115} \\
&\quad{}- 2\,x_{115} - x_{115}^{2} + x_{120} \\
r_{81} &= x_{11} - x_{119} x_{120} - x_{120} \\
r_{82} &= -x_{10} x_{121} - x_{112} + x_{113} \\
r_{83} &= -1 - x_{92} - 4\,x_{92} x_{122} + x_{114} \\
&\quad{}- 4\,x_{115} x_{122} - 4\,x_{122} \\
r_{84} &= 1 + x_{6} + x_{6} x_{92} + x_{6} x_{115} \\
&\quad{}+ x_{92} - x_{111} - x_{111} x_{124} + x_{112} \\
&\quad{}+ x_{115} - x_{116} - x_{123} x_{124} - x_{124} \\
r_{85} &= -3 - 2\,x_{6} - x_{6}^{2} + 2\,x_{6} x_{8} \\
&\quad{}+ 2\,x_{8} - x_{111} - x_{123} \\
r_{86} &= -x_{6} x_{125} - x_{125} + x_{129} \\
r_{87} &= -x_{96} x_{111} + x_{97} - x_{111} x_{129} - x_{126} \\
r_{88} &= -1 + x_{111} - x_{126} - x_{127} \\
r_{89} &= x_{6} - x_{96} - x_{128} \\
r_{90} &= 1 + x_{92} + x_{92} x_{95} + x_{95} \\
&\quad{}+ x_{95} x_{115} + x_{112} + x_{115} - x_{116} \\
&\quad{}- x_{130} - x_{130} x_{132} - x_{131} x_{132} - x_{132} \\
r_{91} &= -3 + 2\,x_{8} + 2\,x_{8} x_{95} - 2\,x_{95} \\
&\quad{}- x_{95}^{2} - x_{130} - x_{131} \\
r_{92} &= -1 - x_{93} x_{96} + x_{96} + x_{130} \\
&\quad{}- x_{133} \\
r_{93} &= -1 + x_{96} - x_{133} - x_{134} \\
r_{94} &= -x_{0} x_{93} - 2\,x_{93} + x_{135} \\
r_{95} &= 1 + x_{0} + x_{0} x_{135} + x_{28} \\
&\quad{}+ x_{31} + x_{31} x_{135} - x_{32} + x_{135} \\
&\quad{}- x_{136} - x_{136} x_{138} - x_{137} x_{138} - x_{138} \\
r_{96} &= -3 + 2\,x_{8} + 2\,x_{8} x_{135} - 2\,x_{135} \\
&\quad{}- x_{135}^{2} - x_{136} - x_{137} \\
r_{97} &= x_{0} x_{93} + x_{28} + x_{31} x_{93} - x_{32} \\
&\quad{}+ x_{93} - x_{139} - x_{139} x_{141} - x_{140} x_{141} \\
&\quad{}- x_{141} \\
r_{98} &= -2 + 2\,x_{8} x_{93} - x_{93}^{2} - x_{139} \\
&\quad{}- x_{140} \\
r_{99} &= -1 - 2\,x_{0} - x_{0}^{2} + x_{143} \\
r_{100} &= -x_{142} x_{143} + x_{144} \\
r_{101} &= -x_{0} x_{8} - x_{8} - x_{8} x_{144} - x_{8} x_{150} \\
&\quad{}+ x_{151} \\
r_{102} &= 2\,x_{0} + x_{0}^{2} + 2\,x_{0} x_{144} + 2\,x_{0} x_{150} \\
&\quad{}+ 2\,x_{144} + x_{144}^{2} + 2\,x_{144} x_{150} + x_{147}^{2} \\
&\quad{}+ 2\,x_{150} + x_{150}^{2} - x_{151}^{2} \\
r_{103} &= -x_{12} x_{149} + x_{152} \\
r_{104} &= -1 + x_{148}^{2} + x_{149}^{2} - x_{152}^{2} \\
r_{105} &= -5 - 4\,x_{0} - 4\,x_{0} x_{153} + x_{12} \\
&\quad{}- 4\,x_{144} - 4\,x_{144} x_{153} - 4\,x_{150} - 4\,x_{150} x_{153} \\
&\quad{}- 4\,x_{153} \\
r_{106} &= -1 - 2\,x_{0} - x_{0}^{2} - 2\,x_{0} x_{144} \\
&\quad{}- 2\,x_{0} x_{150} - 2\,x_{144} - x_{144}^{2} - 2\,x_{144} x_{150} \\
&\quad{}- 2\,x_{150} - x_{150}^{2} + x_{155} \\
r_{107} &= x_{11} - x_{154} x_{155} - x_{155} \\
r_{108} &= -x_{10} x_{156} - x_{147} + x_{148} \\
r_{109} &= -1 - x_{0} - 4\,x_{0} x_{157} - x_{144} \\
&\quad{}- 4\,x_{144} x_{157} + x_{149} - 4\,x_{150} x_{157} - 4\,x_{157} \\
r_{110} &= 2 + 2\,x_{0} + x_{0} x_{6} + x_{6} \\
&\quad{}+ x_{6} x_{144} + x_{6} x_{150} + 2\,x_{144} - x_{146} \\
&\quad{}- x_{146} x_{159} + x_{147} + 2\,x_{150} - x_{151} \\
&\quad{}- x_{158} x_{159} - x_{159} \\
r_{111} &= -6 - 4\,x_{6} - x_{6}^{2} + 2\,x_{6} x_{8} \\
&\quad{}+ 4\,x_{8} - x_{146} - x_{158} \\
r_{112} &= 1 + x_{0} + x_{0} x_{6} + x_{6} \\
&\quad{}+ x_{6} x_{31} + x_{28} + x_{31} - x_{32} \\
&\quad{}- x_{160} - x_{160} x_{162} - x_{161} x_{162} - x_{162} \\
r_{113} &= -3 - 2\,x_{6} - x_{6}^{2} + 2\,x_{6} x_{8} \\
&\quad{}+ 2\,x_{8} - x_{160} - x_{161} \\
r_{114} &= -x_{6} x_{163} - x_{163} + x_{167} \\
r_{115} &= -x_{145} x_{160} + x_{146} - x_{160} x_{167} - x_{164} \\
r_{116} &= -1 + x_{160} - x_{164} - x_{165} \\
r_{117} &= x_{6} - x_{145} - x_{166} \\
r_{118} &= 1 + x_{0} + x_{0} x_{144} + x_{28} \\
&\quad{}+ x_{31} + x_{31} x_{144} - x_{32} + x_{144} \\
&\quad{}- x_{168} - x_{168} x_{170} - x_{169} x_{170} - x_{170} \\
r_{119} &= -3 + 2\,x_{8} + 2\,x_{8} x_{144} - 2\,x_{144} \\
&\quad{}- x_{144}^{2} - x_{168} - x_{169} \\
r_{120} &= -1 - x_{142} x_{145} + x_{145} + x_{168} \\
&\quad{}- x_{171} \\
r_{121} &= -1 + x_{145} - x_{171} - x_{172} \\
r_{122} &= -x_{0} x_{8} - 3\,x_{8} - x_{8} x_{136} - x_{8} x_{178} \\
&\quad{}+ x_{179} \\
r_{123} &= 8 + 6\,x_{0} + x_{0}^{2} + 2\,x_{0} x_{136} \\
&\quad{}+ 2\,x_{0} x_{178} + 6\,x_{136} + x_{136}^{2} + 2\,x_{136} x_{178} \\
&\quad{}+ x_{175}^{2} + 6\,x_{178} + x_{178}^{2} - x_{179}^{2} \\
r_{124} &= -x_{12} x_{177} + x_{180} \\
r_{125} &= -1 + x_{176}^{2} + x_{177}^{2} - x_{180}^{2} \\
r_{126} &= -13 - 4\,x_{0} - 4\,x_{0} x_{181} + x_{12} \\
&\quad{}- 4\,x_{136} - 4\,x_{136} x_{181} - 4\,x_{178} - 4\,x_{178} x_{181} \\
&\quad{}- 12\,x_{181} \\
r_{127} &= -9 - 6\,x_{0} - x_{0}^{2} - 2\,x_{0} x_{136} \\
&\quad{}- 2\,x_{0} x_{178} - 6\,x_{136} - x_{136}^{2} - 2\,x_{136} x_{178} \\
&\quad{}- 6\,x_{178} - x_{178}^{2} + x_{183} \\
r_{128} &= x_{11} - x_{182} x_{183} - x_{183} \\
r_{129} &= -x_{10} x_{184} - x_{175} + x_{176} \\
r_{130} &= -3 - x_{0} - 4\,x_{0} x_{185} - x_{136} \\
&\quad{}- 4\,x_{136} x_{185} + x_{177} - 4\,x_{178} x_{185} - 12\,x_{185} \\
r_{131} &= 6 + 2\,x_{0} + x_{0} x_{6} + 3\,x_{6} \\
&\quad{}+ x_{6} x_{136} + x_{6} x_{178} + 2\,x_{136} - x_{174} \\
&\quad{}- x_{174} x_{187} + x_{175} + 2\,x_{178} - x_{179} \\
&\quad{}- x_{186} x_{187} - x_{187} \\
r_{132} &= -6 - 4\,x_{6} - x_{6}^{2} + 2\,x_{6} x_{8} \\
&\quad{}+ 4\,x_{8} - x_{174} - x_{186} \\
r_{133} &= -x_{6} x_{188} - x_{188} + x_{192} \\
r_{134} &= -x_{160} x_{173} - x_{160} x_{192} + x_{174} - x_{189} \\
r_{135} &= -1 + x_{160} - x_{189} - x_{190} \\
r_{136} &= x_{6} - x_{173} - x_{191} \\
r_{137} &= -x_{139} x_{142} + x_{193} \\
r_{138} &= 1 - 2\,x_{93} - x_{93} x_{136} + x_{194} \\
r_{139} &= x_{173} x_{193} - x_{194} x_{196} - x_{195} \\
r_{140} &= -1 + x_{194} - x_{195} - x_{197} \\
r_{141} &= -x_{8} - x_{8} x_{47} - x_{8} x_{203} + x_{204} \\
r_{142} &= 2\,x_{47} + x_{47}^{2} + 2\,x_{47} x_{203} + x_{200}^{2} \\
&\quad{}+ 2\,x_{203} + x_{203}^{2} - x_{204}^{2} \\
r_{143} &= -x_{12} x_{202} + x_{205} \\
r_{144} &= -1 + x_{201}^{2} + x_{202}^{2} - x_{205}^{2} \\
r_{145} &= -5 + x_{12} - 4\,x_{47} - 4\,x_{47} x_{206} \\
&\quad{}- 4\,x_{203} - 4\,x_{203} x_{206} - 4\,x_{206} \\
r_{146} &= -1 - 2\,x_{47} - x_{47}^{2} - 2\,x_{47} x_{203} \\
&\quad{}- 2\,x_{203} - x_{203}^{2} + x_{208} \\
r_{147} &= x_{11} - x_{207} x_{208} - x_{208} \\
r_{148} &= -x_{10} x_{209} - x_{200} + x_{201} \\
r_{149} &= -1 - x_{47} - 4\,x_{47} x_{210} + x_{202} \\
&\quad{}- 4\,x_{203} x_{210} - 4\,x_{210} \\
r_{150} &= 2 + x_{6} + x_{6} x_{47} + x_{6} x_{203} \\
&\quad{}+ 2\,x_{47} - x_{199} - x_{199} x_{212} + x_{200} \\
&\quad{}+ 2\,x_{203} - x_{204} - x_{211} x_{212} - x_{212} \\
r_{151} &= -6 - 4\,x_{6} - x_{6}^{2} + 2\,x_{6} x_{8} \\
&\quad{}+ 4\,x_{8} - x_{199} - x_{211} \\
r_{152} &= -x_{3} x_{8} - 2\,x_{8} - x_{8} x_{217} + x_{218} \\
r_{153} &= 3 + 4\,x_{3} + x_{3}^{2} + 2\,x_{3} x_{217} \\
&\quad{}+ x_{214}^{2} + 4\,x_{217} + x_{217}^{2} - x_{218}^{2} \\
r_{154} &= -x_{12} x_{216} + x_{219} \\
r_{155} &= -1 + x_{215}^{2} + x_{216}^{2} - x_{219}^{2} \\
r_{156} &= -9 - 4\,x_{3} - 4\,x_{3} x_{220} + x_{12} \\
&\quad{}- 4\,x_{217} - 4\,x_{217} x_{220} - 8\,x_{220} \\
r_{157} &= -4 - 4\,x_{3} - x_{3}^{2} - 2\,x_{3} x_{217} \\
&\quad{}- 4\,x_{217} - x_{217}^{2} + x_{222} \\
r_{158} &= x_{11} - x_{221} x_{222} - x_{222} \\
r_{159} &= -x_{10} x_{223} - x_{214} + x_{215} \\
r_{160} &= -2 - x_{3} - 4\,x_{3} x_{224} + x_{216} \\
&\quad{}- 4\,x_{217} x_{224} - 8\,x_{224} \\
r_{161} &= 2 + x_{3} + x_{3} x_{6} + 2\,x_{6} \\
&\quad{}+ x_{6} x_{217} - x_{213} - x_{213} x_{226} + x_{214} \\
&\quad{}+ x_{217} - x_{218} - x_{225} x_{226} - x_{226} \\
r_{162} &= -3 - 2\,x_{6} - x_{6}^{2} + 2\,x_{6} x_{8} \\
&\quad{}+ 2\,x_{8} - x_{213} - x_{225} \\
r_{163} &= -x_{6} x_{227} - x_{227} + x_{231} \\
r_{164} &= -x_{198} x_{213} + x_{199} - x_{213} x_{231} - x_{228} \\
r_{165} &= -1 + x_{213} - x_{228} - x_{229} \\
r_{166} &= x_{6} - x_{198} - x_{230} \\
r_{167} &= -x_{195} x_{232} + x_{198} \\
r_{168} &= -x_{8} - x_{8} x_{48} - x_{8} x_{238} + x_{239} \\
r_{169} &= 2\,x_{48} + x_{48}^{2} + 2\,x_{48} x_{238} + x_{235}^{2} \\
&\quad{}+ 2\,x_{238} + x_{238}^{2} - x_{239}^{2} \\
r_{170} &= -x_{12} x_{237} + x_{240} \\
r_{171} &= -1 + x_{236}^{2} + x_{237}^{2} - x_{240}^{2} \\
r_{172} &= -5 + x_{12} - 4\,x_{48} - 4\,x_{48} x_{241} \\
&\quad{}- 4\,x_{238} - 4\,x_{238} x_{241} - 4\,x_{241} \\
r_{173} &= -1 - 2\,x_{48} - x_{48}^{2} - 2\,x_{48} x_{238} \\
&\quad{}- 2\,x_{238} - x_{238}^{2} + x_{243} \\
r_{174} &= x_{11} - x_{242} x_{243} - x_{243} \\
r_{175} &= -x_{10} x_{244} - x_{235} + x_{236} \\
r_{176} &= -1 - x_{48} - 4\,x_{48} x_{245} + x_{237} \\
&\quad{}- 4\,x_{238} x_{245} - 4\,x_{245} \\
r_{177} &= 2 + x_{6} + x_{6} x_{48} + x_{6} x_{238} \\
&\quad{}+ 2\,x_{48} - x_{234} - x_{234} x_{247} + x_{235} \\
&\quad{}+ 2\,x_{238} - x_{239} - x_{246} x_{247} - x_{247} \\
r_{178} &= -6 - 4\,x_{6} - x_{6}^{2} + 2\,x_{6} x_{8} \\
&\quad{}+ 4\,x_{8} - x_{234} - x_{246} \\
r_{179} &= -x_{6} x_{248} - x_{248} + x_{252} \\
r_{180} &= -x_{213} x_{233} - x_{213} x_{252} + x_{234} - x_{249} \\
r_{181} &= -1 + x_{213} - x_{249} - x_{250} \\
r_{182} &= x_{6} - x_{233} - x_{251} \\
r_{183} &= -x_{195} x_{253} + x_{233} \\
r_{184} &= -x_{8} - x_{8} x_{49} - x_{8} x_{259} + x_{260} \\
r_{185} &= 2\,x_{49} + x_{49}^{2} + 2\,x_{49} x_{259} + x_{256}^{2} \\
&\quad{}+ 2\,x_{259} + x_{259}^{2} - x_{260}^{2} \\
r_{186} &= -x_{12} x_{258} + x_{261} \\
r_{187} &= -1 + x_{257}^{2} + x_{258}^{2} - x_{261}^{2} \\
r_{188} &= -5 + x_{12} - 4\,x_{49} - 4\,x_{49} x_{262} \\
&\quad{}- 4\,x_{259} - 4\,x_{259} x_{262} - 4\,x_{262} \\
r_{189} &= -1 - 2\,x_{49} - x_{49}^{2} - 2\,x_{49} x_{259} \\
&\quad{}- 2\,x_{259} - x_{259}^{2} + x_{264} \\
r_{190} &= x_{11} - x_{263} x_{264} - x_{264} \\
r_{191} &= -x_{10} x_{265} - x_{256} + x_{257} \\
r_{192} &= -1 - x_{49} - 4\,x_{49} x_{266} + x_{258} \\
&\quad{}- 4\,x_{259} x_{266} - 4\,x_{266} \\
r_{193} &= 2 + x_{6} + x_{6} x_{49} + x_{6} x_{259} \\
&\quad{}+ 2\,x_{49} - x_{255} - x_{255} x_{268} + x_{256} \\
&\quad{}+ 2\,x_{259} - x_{260} - x_{267} x_{268} - x_{268} \\
r_{194} &= -6 - 4\,x_{6} - x_{6}^{2} + 2\,x_{6} x_{8} \\
&\quad{}+ 4\,x_{8} - x_{255} - x_{267} \\
r_{195} &= -x_{6} x_{269} - x_{269} + x_{273} \\
r_{196} &= -x_{213} x_{254} - x_{213} x_{273} + x_{255} - x_{270} \\
r_{197} &= -1 + x_{213} - x_{270} - x_{271} \\
r_{198} &= x_{6} - x_{254} - x_{272} \\
r_{199} &= -x_{195} x_{274} + x_{254} \\
r_{200} &= -x_{8} - x_{8} x_{50} - x_{8} x_{280} + x_{281} \\
r_{201} &= 2\,x_{50} + x_{50}^{2} + 2\,x_{50} x_{280} + x_{277}^{2} \\
&\quad{}+ 2\,x_{280} + x_{280}^{2} - x_{281}^{2} \\
r_{202} &= -x_{12} x_{279} + x_{282} \\
r_{203} &= -1 + x_{278}^{2} + x_{279}^{2} - x_{282}^{2} \\
r_{204} &= -5 + x_{12} - 4\,x_{50} - 4\,x_{50} x_{283} \\
&\quad{}- 4\,x_{280} - 4\,x_{280} x_{283} - 4\,x_{283} \\
r_{205} &= -1 - 2\,x_{50} - x_{50}^{2} - 2\,x_{50} x_{280} \\
&\quad{}- 2\,x_{280} - x_{280}^{2} + x_{285} \\
r_{206} &= x_{11} - x_{284} x_{285} - x_{285} \\
r_{207} &= -x_{10} x_{286} - x_{277} + x_{278} \\
r_{208} &= -1 - x_{50} - 4\,x_{50} x_{287} + x_{279} \\
&\quad{}- 4\,x_{280} x_{287} - 4\,x_{287} \\
r_{209} &= 2 + x_{6} + x_{6} x_{50} + x_{6} x_{280} \\
&\quad{}+ 2\,x_{50} - x_{276} - x_{276} x_{289} + x_{277} \\
&\quad{}+ 2\,x_{280} - x_{281} - x_{288} x_{289} - x_{289} \\
r_{210} &= -6 - 4\,x_{6} - x_{6}^{2} + 2\,x_{6} x_{8} \\
&\quad{}+ 4\,x_{8} - x_{276} - x_{288} \\
r_{211} &= -2 - x_{2} - x_{3} - x_{46} \\
&\quad{}- x_{294} + x_{295} \\
r_{212} &= -x_{8} x_{295} + x_{296} \\
r_{213} &= -1 + x_{291}^{2} + x_{295}^{2} - x_{296}^{2} \\
r_{214} &= -x_{12} x_{293} + x_{297} \\
r_{215} &= -1 + x_{292}^{2} + x_{293}^{2} - x_{297}^{2} \\
r_{216} &= -1 + x_{12} - 4\,x_{295} - 4\,x_{295} x_{298} \\
r_{217} &= -x_{295}^{2} + x_{300} \\
r_{218} &= x_{11} - x_{299} x_{300} - x_{300} \\
r_{219} &= -x_{10} x_{301} - x_{291} + x_{292} \\
r_{220} &= -2 - x_{2} - x_{3} - x_{46} \\
&\quad{}+ x_{293} - 4\,x_{295} x_{302} \\
r_{221} &= x_{6} x_{295} - x_{290} - x_{290} x_{304} + x_{291} \\
&\quad{}+ x_{295} - x_{296} - x_{303} x_{304} - x_{304} \\
r_{222} &= -3 - 2\,x_{6} - x_{6}^{2} + 2\,x_{6} x_{8} \\
&\quad{}+ 2\,x_{8} - x_{290} - x_{303} \\
r_{223} &= -x_{6} x_{305} - x_{305} + x_{309} \\
r_{224} &= -x_{275} x_{290} + x_{276} - x_{290} x_{309} - x_{306} \\
r_{225} &= -1 + x_{290} - x_{306} - x_{307} \\
r_{226} &= x_{6} - x_{275} - x_{308} \\
r_{227} &= -x_{195} x_{310} + x_{275} \\
r_{228} &= -x_{8} - x_{8} x_{51} - x_{8} x_{316} + x_{317} \\
r_{229} &= 2\,x_{51} + x_{51}^{2} + 2\,x_{51} x_{316} + x_{313}^{2} \\
&\quad{}+ 2\,x_{316} + x_{316}^{2} - x_{317}^{2} \\
r_{230} &= -x_{12} x_{315} + x_{318} \\
r_{231} &= -1 + x_{314}^{2} + x_{315}^{2} - x_{318}^{2} \\
r_{232} &= -5 + x_{12} - 4\,x_{51} - 4\,x_{51} x_{319} \\
&\quad{}- 4\,x_{316} - 4\,x_{316} x_{319} - 4\,x_{319} \\
r_{233} &= -1 - 2\,x_{51} - x_{51}^{2} - 2\,x_{51} x_{316} \\
&\quad{}- 2\,x_{316} - x_{316}^{2} + x_{321} \\
r_{234} &= x_{11} - x_{320} x_{321} - x_{321} \\
r_{235} &= -x_{10} x_{322} - x_{313} + x_{314} \\
r_{236} &= -1 - x_{51} - 4\,x_{51} x_{323} + x_{315} \\
&\quad{}- 4\,x_{316} x_{323} - 4\,x_{323} \\
r_{237} &= 2 + x_{6} + x_{6} x_{51} + x_{6} x_{316} \\
&\quad{}+ 2\,x_{51} - x_{312} - x_{312} x_{325} + x_{313} \\
&\quad{}+ 2\,x_{316} - x_{317} - x_{324} x_{325} - x_{325} \\
r_{238} &= -6 - 4\,x_{6} - x_{6}^{2} + 2\,x_{6} x_{8} \\
&\quad{}+ 4\,x_{8} - x_{312} - x_{324} \\
r_{239} &= -x_{6} x_{326} - x_{326} + x_{330} \\
r_{240} &= -x_{290} x_{311} - x_{290} x_{330} + x_{312} - x_{327} \\
r_{241} &= -1 + x_{290} - x_{327} - x_{328} \\
r_{242} &= x_{6} - x_{311} - x_{329} \\
r_{243} &= -x_{195} x_{331} + x_{311} \\
r_{244} &= -x_{8} - x_{8} x_{52} - x_{8} x_{337} + x_{338} \\
r_{245} &= 2\,x_{52} + x_{52}^{2} + 2\,x_{52} x_{337} + x_{334}^{2} \\
&\quad{}+ 2\,x_{337} + x_{337}^{2} - x_{338}^{2} \\
r_{246} &= -x_{12} x_{336} + x_{339} \\
r_{247} &= -1 + x_{335}^{2} + x_{336}^{2} - x_{339}^{2} \\
r_{248} &= -5 + x_{12} - 4\,x_{52} - 4\,x_{52} x_{340} \\
&\quad{}- 4\,x_{337} - 4\,x_{337} x_{340} - 4\,x_{340} \\
r_{249} &= -1 - 2\,x_{52} - x_{52}^{2} - 2\,x_{52} x_{337} \\
&\quad{}- 2\,x_{337} - x_{337}^{2} + x_{342} \\
r_{250} &= x_{11} - x_{341} x_{342} - x_{342} \\
r_{251} &= -x_{10} x_{343} - x_{334} + x_{335} \\
r_{252} &= -1 - x_{52} - 4\,x_{52} x_{344} + x_{336} \\
&\quad{}- 4\,x_{337} x_{344} - 4\,x_{344} \\
r_{253} &= 2 + x_{6} + x_{6} x_{52} + x_{6} x_{337} \\
&\quad{}+ 2\,x_{52} - x_{333} - x_{333} x_{346} + x_{334} \\
&\quad{}+ 2\,x_{337} - x_{338} - x_{345} x_{346} - x_{346} \\
r_{254} &= -6 - 4\,x_{6} - x_{6}^{2} + 2\,x_{6} x_{8} \\
&\quad{}+ 4\,x_{8} - x_{333} - x_{345} \\
r_{255} &= -x_{6} x_{347} - x_{347} + x_{351} \\
r_{256} &= -x_{290} x_{332} - x_{290} x_{351} + x_{333} - x_{348} \\
r_{257} &= -1 + x_{290} - x_{348} - x_{349} \\
r_{258} &= x_{6} - x_{332} - x_{350} \\
r_{259} &= -x_{195} x_{352} + x_{332} \\
r_{260} &= -x_{8} - x_{8} x_{53} - x_{8} x_{358} + x_{359} \\
r_{261} &= 2\,x_{53} + x_{53}^{2} + 2\,x_{53} x_{358} + x_{355}^{2} \\
&\quad{}+ 2\,x_{358} + x_{358}^{2} - x_{359}^{2} \\
r_{262} &= -x_{12} x_{357} + x_{360} \\
r_{263} &= -1 + x_{356}^{2} + x_{357}^{2} - x_{360}^{2} \\
r_{264} &= -5 + x_{12} - 4\,x_{53} - 4\,x_{53} x_{361} \\
&\quad{}- 4\,x_{358} - 4\,x_{358} x_{361} - 4\,x_{361} \\
r_{265} &= -1 - 2\,x_{53} - x_{53}^{2} - 2\,x_{53} x_{358} \\
&\quad{}- 2\,x_{358} - x_{358}^{2} + x_{363} \\
r_{266} &= x_{11} - x_{362} x_{363} - x_{363} \\
r_{267} &= -x_{10} x_{364} - x_{355} + x_{356} \\
r_{268} &= -1 - x_{53} - 4\,x_{53} x_{365} + x_{357} \\
&\quad{}- 4\,x_{358} x_{365} - 4\,x_{365} \\
r_{269} &= 2 + x_{6} + x_{6} x_{53} + x_{6} x_{358} \\
&\quad{}+ 2\,x_{53} - x_{354} - x_{354} x_{367} + x_{355} \\
&\quad{}+ 2\,x_{358} - x_{359} - x_{366} x_{367} - x_{367} \\
r_{270} &= -6 - 4\,x_{6} - x_{6}^{2} + 2\,x_{6} x_{8} \\
&\quad{}+ 4\,x_{8} - x_{354} - x_{366} \\
r_{271} &= -x_{6} x_{368} - x_{368} + x_{372} \\
r_{272} &= -x_{290} x_{353} - x_{290} x_{372} + x_{354} - x_{369} \\
r_{273} &= -1 + x_{290} - x_{369} - x_{370} \\
r_{274} &= x_{6} - x_{353} - x_{371} \\
r_{275} &= -x_{195} x_{373} + x_{353} \\
r_{276} &= -x_{8} - x_{8} x_{54} - x_{8} x_{379} + x_{380} \\
r_{277} &= 2\,x_{54} + x_{54}^{2} + 2\,x_{54} x_{379} + x_{376}^{2} \\
&\quad{}+ 2\,x_{379} + x_{379}^{2} - x_{380}^{2} \\
r_{278} &= -x_{12} x_{378} + x_{381} \\
r_{279} &= -1 + x_{377}^{2} + x_{378}^{2} - x_{381}^{2} \\
r_{280} &= -5 + x_{12} - 4\,x_{54} - 4\,x_{54} x_{382} \\
&\quad{}- 4\,x_{379} - 4\,x_{379} x_{382} - 4\,x_{382} \\
r_{281} &= -1 - 2\,x_{54} - x_{54}^{2} - 2\,x_{54} x_{379} \\
&\quad{}- 2\,x_{379} - x_{379}^{2} + x_{384} \\
r_{282} &= x_{11} - x_{383} x_{384} - x_{384} \\
r_{283} &= -x_{10} x_{385} - x_{376} + x_{377} \\
r_{284} &= -1 - x_{54} - 4\,x_{54} x_{386} + x_{378} \\
&\quad{}- 4\,x_{379} x_{386} - 4\,x_{386} \\
r_{285} &= 2 + x_{6} + x_{6} x_{54} + x_{6} x_{379} \\
&\quad{}+ 2\,x_{54} - x_{375} - x_{375} x_{388} + x_{376} \\
&\quad{}+ 2\,x_{379} - x_{380} - x_{387} x_{388} - x_{388} \\
r_{286} &= -6 - 4\,x_{6} - x_{6}^{2} + 2\,x_{6} x_{8} \\
&\quad{}+ 4\,x_{8} - x_{375} - x_{387} \\
r_{287} &= -x_{2} x_{8} - x_{8} - x_{8} x_{393} + x_{394} \\
r_{288} &= 2\,x_{2} + x_{2}^{2} + 2\,x_{2} x_{393} + x_{390}^{2} \\
&\quad{}+ 2\,x_{393} + x_{393}^{2} - x_{394}^{2} \\
r_{289} &= -x_{12} x_{392} + x_{395} \\
r_{290} &= -1 + x_{391}^{2} + x_{392}^{2} - x_{395}^{2} \\
r_{291} &= -5 - 4\,x_{2} - 4\,x_{2} x_{396} + x_{12} \\
&\quad{}- 4\,x_{393} - 4\,x_{393} x_{396} - 4\,x_{396} \\
r_{292} &= -1 - 2\,x_{2} - x_{2}^{2} - 2\,x_{2} x_{393} \\
&\quad{}- 2\,x_{393} - x_{393}^{2} + x_{398} \\
r_{293} &= x_{11} - x_{397} x_{398} - x_{398} \\
r_{294} &= -x_{10} x_{399} - x_{390} + x_{391} \\
r_{295} &= -1 - x_{2} - 4\,x_{2} x_{400} + x_{392} \\
&\quad{}- 4\,x_{393} x_{400} - 4\,x_{400} \\
r_{296} &= 1 + x_{2} + x_{2} x_{6} + x_{6} \\
&\quad{}+ x_{6} x_{393} - x_{389} - x_{389} x_{402} + x_{390} \\
&\quad{}+ x_{393} - x_{394} - x_{401} x_{402} - x_{402} \\
r_{297} &= -3 - 2\,x_{6} - x_{6}^{2} + 2\,x_{6} x_{8} \\
&\quad{}+ 2\,x_{8} - x_{389} - x_{401} \\
r_{298} &= -x_{6} x_{403} - x_{403} + x_{407} \\
r_{299} &= -x_{374} x_{389} + x_{375} - x_{389} x_{407} - x_{404} \\
r_{300} &= -1 + x_{389} - x_{404} - x_{405} \\
r_{301} &= x_{6} - x_{374} - x_{406} \\
r_{302} &= -x_{195} x_{408} + x_{374} \\
r_{303} &= -x_{8} - x_{8} x_{55} - x_{8} x_{414} + x_{415} \\
r_{304} &= 2\,x_{55} + x_{55}^{2} + 2\,x_{55} x_{414} + x_{411}^{2} \\
&\quad{}+ 2\,x_{414} + x_{414}^{2} - x_{415}^{2} \\
r_{305} &= -x_{12} x_{413} + x_{416} \\
r_{306} &= -1 + x_{412}^{2} + x_{413}^{2} - x_{416}^{2} \\
r_{307} &= -5 + x_{12} - 4\,x_{55} - 4\,x_{55} x_{417} \\
&\quad{}- 4\,x_{414} - 4\,x_{414} x_{417} - 4\,x_{417} \\
r_{308} &= -1 - 2\,x_{55} - x_{55}^{2} - 2\,x_{55} x_{414} \\
&\quad{}- 2\,x_{414} - x_{414}^{2} + x_{419} \\
r_{309} &= x_{11} - x_{418} x_{419} - x_{419} \\
r_{310} &= -x_{10} x_{420} - x_{411} + x_{412} \\
r_{311} &= -1 - x_{55} - 4\,x_{55} x_{421} + x_{413} \\
&\quad{}- 4\,x_{414} x_{421} - 4\,x_{421} \\
r_{312} &= 2 + x_{6} + x_{6} x_{55} + x_{6} x_{414} \\
&\quad{}+ 2\,x_{55} - x_{410} - x_{410} x_{423} + x_{411} \\
&\quad{}+ 2\,x_{414} - x_{415} - x_{422} x_{423} - x_{423} \\
r_{313} &= -6 - 4\,x_{6} - x_{6}^{2} + 2\,x_{6} x_{8} \\
&\quad{}+ 4\,x_{8} - x_{410} - x_{422} \\
r_{314} &= -x_{6} x_{424} - x_{424} + x_{428} \\
r_{315} &= -x_{389} x_{409} - x_{389} x_{428} + x_{410} - x_{425} \\
r_{316} &= -1 + x_{389} - x_{425} - x_{426} \\
r_{317} &= x_{6} - x_{409} - x_{427} \\
r_{318} &= -x_{195} x_{429} + x_{409} \\
r_{319} &= -x_{8} - x_{8} x_{56} - x_{8} x_{435} + x_{436} \\
r_{320} &= 2\,x_{56} + x_{56}^{2} + 2\,x_{56} x_{435} + x_{432}^{2} \\
&\quad{}+ 2\,x_{435} + x_{435}^{2} - x_{436}^{2} \\
r_{321} &= -x_{12} x_{434} + x_{437} \\
r_{322} &= -1 + x_{433}^{2} + x_{434}^{2} - x_{437}^{2} \\
r_{323} &= -5 + x_{12} - 4\,x_{56} - 4\,x_{56} x_{438} \\
&\quad{}- 4\,x_{435} - 4\,x_{435} x_{438} - 4\,x_{438} \\
r_{324} &= -1 - 2\,x_{56} - x_{56}^{2} - 2\,x_{56} x_{435} \\
&\quad{}- 2\,x_{435} - x_{435}^{2} + x_{440} \\
r_{325} &= x_{11} - x_{439} x_{440} - x_{440} \\
r_{326} &= -x_{10} x_{441} - x_{432} + x_{433} \\
r_{327} &= -1 - x_{56} - 4\,x_{56} x_{442} + x_{434} \\
&\quad{}- 4\,x_{435} x_{442} - 4\,x_{442} \\
r_{328} &= 2 + x_{6} + x_{6} x_{56} + x_{6} x_{435} \\
&\quad{}+ 2\,x_{56} - x_{431} - x_{431} x_{444} + x_{432} \\
&\quad{}+ 2\,x_{435} - x_{436} - x_{443} x_{444} - x_{444} \\
r_{329} &= -6 - 4\,x_{6} - x_{6}^{2} + 2\,x_{6} x_{8} \\
&\quad{}+ 4\,x_{8} - x_{431} - x_{443} \\
r_{330} &= -x_{6} x_{445} - x_{445} + x_{449} \\
r_{331} &= -x_{290} x_{430} - x_{290} x_{449} + x_{431} - x_{446} \\
r_{332} &= -1 + x_{290} - x_{446} - x_{447} \\
r_{333} &= x_{6} - x_{430} - x_{448} \\
r_{334} &= -x_{195} x_{450} + x_{430} \\
r_{335} &= -x_{8} - x_{8} x_{57} - x_{8} x_{456} + x_{457} \\
r_{336} &= 2\,x_{57} + x_{57}^{2} + 2\,x_{57} x_{456} + x_{453}^{2} \\
&\quad{}+ 2\,x_{456} + x_{456}^{2} - x_{457}^{2} \\
r_{337} &= -x_{12} x_{455} + x_{458} \\
r_{338} &= -1 + x_{454}^{2} + x_{455}^{2} - x_{458}^{2} \\
r_{339} &= -5 + x_{12} - 4\,x_{57} - 4\,x_{57} x_{459} \\
&\quad{}- 4\,x_{456} - 4\,x_{456} x_{459} - 4\,x_{459} \\
r_{340} &= -1 - 2\,x_{57} - x_{57}^{2} - 2\,x_{57} x_{456} \\
&\quad{}- 2\,x_{456} - x_{456}^{2} + x_{461} \\
r_{341} &= x_{11} - x_{460} x_{461} - x_{461} \\
r_{342} &= -x_{10} x_{462} - x_{453} + x_{454} \\
r_{343} &= -1 - x_{57} - 4\,x_{57} x_{463} + x_{455} \\
&\quad{}- 4\,x_{456} x_{463} - 4\,x_{463} \\
r_{344} &= 2 + x_{6} + x_{6} x_{57} + x_{6} x_{456} \\
&\quad{}+ 2\,x_{57} - x_{452} - x_{452} x_{465} + x_{453} \\
&\quad{}+ 2\,x_{456} - x_{457} - x_{464} x_{465} - x_{465} \\
r_{345} &= -6 - 4\,x_{6} - x_{6}^{2} + 2\,x_{6} x_{8} \\
&\quad{}+ 4\,x_{8} - x_{452} - x_{464} \\
r_{346} &= -x_{6} x_{466} - x_{466} + x_{470} \\
r_{347} &= -x_{290} x_{451} - x_{290} x_{470} + x_{452} - x_{467} \\
r_{348} &= -1 + x_{290} - x_{467} - x_{468} \\
r_{349} &= x_{6} - x_{451} - x_{469} \\
r_{350} &= -x_{195} x_{471} + x_{451} \\
r_{351} &= 1 + x_{58} x_{93} + x_{93} - x_{195} x_{472} \\
r_{352} &= -6 - 5\,x_{50} - x_{50}^{2} - 2\,x_{50} x_{51} \\
&\quad{}- 5\,x_{51} - x_{51}^{2} + x_{473} \\
r_{353} &= -72 - 68\,x_{51} - 16\,x_{51}^{2} - 16\,x_{51} x_{473} \\
&\quad{}- 34\,x_{473} - 4\,x_{473}^{2} + x_{474} \\
r_{354} &= -30 - 11\,x_{50} - x_{50}^{2} - 4\,x_{50} x_{474} \\
&\quad{}- 22\,x_{474} - 4\,x_{474}^{2} + x_{475} \\
r_{355} &= -380 - 312\,x_{474} - 64\,x_{474}^{2} - 32\,x_{474} x_{475} \\
&\quad{}- 78\,x_{475} - 4\,x_{475}^{2} + x_{476} \\
r_{356} &= -6 - 5\,x_{51} - x_{51}^{2} - 2\,x_{51} x_{52} \\
&\quad{}- 5\,x_{52} - x_{52}^{2} + x_{477} \\
r_{357} &= -6 - 5\,x_{50} - x_{50}^{2} - 2\,x_{50} x_{52} \\
&\quad{}- 5\,x_{52} - x_{52}^{2} + x_{478} \\
r_{358} &= -12 - 4\,x_{51} - 4\,x_{52} - 2\,x_{473} \\
&\quad{}- 2\,x_{478} + x_{479} \\
r_{359} &= -102 - 22\,x_{50} - 2\,x_{50}^{2} - 16\,x_{50} x_{52} \\
&\quad{}- 8\,x_{50} x_{477} - 120\,x_{52} - 32\,x_{52}^{2} - 32\,x_{52} x_{477} \\
&\quad{}- 52\,x_{477} - 8\,x_{477}^{2} - 8\,x_{478} - 2\,x_{479} \\
&\quad{}- 2\,x_{479}^{2} + x_{480} \\
r_{360} &= -96 - 64\,x_{52} - 32\,x_{478} - 8\,x_{479} \\
&\quad{}- 8\,x_{479}^{2} - 2\,x_{480} - 2\,x_{480}^{2} + x_{481} \\
r_{361} &= -10 - 4\,x_{50} - 4\,x_{51} - 2\,x_{477} \\
&\quad{}- 2\,x_{478} + x_{482} \\
r_{362} &= -27 - 20\,x_{51} - 2\,x_{473} - 8\,x_{477} \\
&\quad{}- 2\,x_{482} - 2\,x_{482}^{2} + x_{483} \\
r_{363} &= -80 - 64\,x_{51} - 32\,x_{477} - 8\,x_{482} \\
&\quad{}- 8\,x_{482}^{2} - 2\,x_{483} - 2\,x_{483}^{2} + x_{484} \\
r_{364} &= -276 - 30\,x_{50} - 2\,x_{50}^{2} - 16\,x_{50} x_{52} \\
&\quad{}- 8\,x_{50} x_{477} - 30\,x_{51} - 2\,x_{51}^{2} - 16\,x_{51} x_{52} \\
&\quad{}- 8\,x_{51} x_{478} - 272\,x_{52} - 64\,x_{52}^{2} - 32\,x_{52} x_{477} \\
&\quad{}- 32\,x_{52} x_{478} - 68\,x_{477} - 8\,x_{477}^{2} - 68\,x_{478} \\
&\quad{}- 8\,x_{478}^{2} + x_{485} \\
r_{365} &= -552 - 120\,x_{51} - 8\,x_{51}^{2} - 64\,x_{51} x_{52} \\
&\quad{}- 32\,x_{51} x_{478} - 544\,x_{52} - 128\,x_{52}^{2} - 128\,x_{52} x_{478} \\
&\quad{}- 272\,x_{478} - 32\,x_{478}^{2} - 2\,x_{485} - 2\,x_{485}^{2} \\
&\quad{}+ x_{486} \\
r_{366} &= -247 - 4\,x_{50} - 352\,x_{51} - 128\,x_{51}^{2} \\
&\quad{}- 64\,x_{51} x_{473} - 64\,x_{51} x_{477} - 92\,x_{473} - 8\,x_{473}^{2} \\
&\quad{}- 16\,x_{473} x_{477} - 84\,x_{477} - 8\,x_{477}^{2} - 2\,x_{478} \\
&\quad{}+ x_{487} \\
r_{367} &= -960 - 1408\,x_{51} - 512\,x_{51}^{2} - 256\,x_{51} x_{473} \\
&\quad{}- 256\,x_{51} x_{477} - 368\,x_{473} - 32\,x_{473}^{2} - 64\,x_{473} x_{477} \\
&\quad{}- 336\,x_{477} - 32\,x_{477}^{2} - 2\,x_{487} - 2\,x_{487}^{2} \\
&\quad{}+ x_{488} \\
r_{368} &= -56 - 60\,x_{51} - 16\,x_{51}^{2} - 16\,x_{51} x_{473} \\
&\quad{}- 30\,x_{473} - 4\,x_{473}^{2} + x_{489} \\
r_{369} &= -30 - 11\,x_{50} - x_{50}^{2} - 4\,x_{50} x_{489} \\
&\quad{}- 22\,x_{489} - 4\,x_{489}^{2} + x_{490} \\
r_{370} &= -380 - 312\,x_{489} - 64\,x_{489}^{2} - 32\,x_{489} x_{490} \\
&\quad{}- 78\,x_{490} - 4\,x_{490}^{2} + x_{491} \\
r_{371} &= -30 - 11\,x_{51} - x_{51}^{2} - 4\,x_{51} x_{491} \\
&\quad{}- 22\,x_{491} - 4\,x_{491}^{2} + x_{492} \\
r_{372} &= -349 - 420\,x_{51} - 128\,x_{51}^{2} - 128\,x_{51} x_{473} \\
&\quad{}- 208\,x_{473} - 32\,x_{473}^{2} - 2\,x_{477} + x_{493} \\
r_{373} &= -1368 - 1664\,x_{51} - 512\,x_{51}^{2} - 512\,x_{51} x_{473} \\
&\quad{}- 832\,x_{473} - 128\,x_{473}^{2} - 2\,x_{493} - 2\,x_{493}^{2} \\
&\quad{}+ x_{494} \\
r_{374} &= -56 - 60\,x_{50} - 16\,x_{50}^{2} - 16\,x_{50} x_{478} \\
&\quad{}- 8\,x_{50} x_{494} - 30\,x_{478} - 4\,x_{478}^{2} - 4\,x_{478} x_{494} \\
&\quad{}- 15\,x_{494} - x_{494}^{2} + x_{495} \\
r_{375} &= -6 - 5\,x_{52} - x_{52}^{2} - 2\,x_{52} x_{53} \\
&\quad{}- 5\,x_{53} - x_{53}^{2} + x_{496} \\
r_{376} &= -224 - 26\,x_{50} - 2\,x_{50}^{2} - 16\,x_{50} x_{53} \\
&\quad{}- 8\,x_{50} x_{496} - 288\,x_{52} - 128\,x_{52}^{2} - 64\,x_{52} x_{477} \\
&\quad{}- 64\,x_{52} x_{478} - 88\,x_{53} - 32\,x_{53}^{2} - 32\,x_{53} x_{496} \\
&\quad{}- 76\,x_{477} - 8\,x_{477}^{2} - 16\,x_{477} x_{478} - 68\,x_{478} \\
&\quad{}- 8\,x_{478}^{2} - 44\,x_{496} - 8\,x_{496}^{2} + x_{497} \\
r_{377} &= -640 - 1152\,x_{52} - 512\,x_{52}^{2} - 256\,x_{52} x_{477} \\
&\quad{}- 256\,x_{52} x_{478} - 304\,x_{477} - 32\,x_{477}^{2} - 64\,x_{477} x_{478} \\
&\quad{}- 272\,x_{478} - 32\,x_{478}^{2} - 2\,x_{497} - 2\,x_{497}^{2} \\
&\quad{}+ x_{498} \\
r_{378} &= -21 - 16\,x_{51} - x_{52} - 8\,x_{477} \\
&\quad{}- 2\,x_{482} - 2\,x_{482}^{2} + x_{499} \\
r_{379} &= -42 - 52\,x_{52} - 16\,x_{52}^{2} - 16\,x_{52} x_{477} \\
&\quad{}- 26\,x_{477} - 4\,x_{477}^{2} + x_{500} \\
r_{380} &= -30 - 11\,x_{52} - x_{52}^{2} - 4\,x_{52} x_{500} \\
&\quad{}- 22\,x_{500} - 4\,x_{500}^{2} + x_{501} \\
r_{381} &= -306 - 35\,x_{51} - x_{51}^{2} - 16\,x_{51} x_{500} \\
&\quad{}- 4\,x_{51} x_{501} - 280\,x_{500} - 64\,x_{500}^{2} - 32\,x_{500} x_{501} \\
&\quad{}- 70\,x_{501} - 4\,x_{501}^{2} + x_{502} \\
r_{382} &= -4556 - 4320\,x_{500} - 1024\,x_{500}^{2} - 512\,x_{500} x_{501} \\
&\quad{}- 128\,x_{500} x_{502} - 1080\,x_{501} - 64\,x_{501}^{2} - 32\,x_{501} x_{502} \\
&\quad{}- 270\,x_{502} - 4\,x_{502}^{2} + x_{503} \\
r_{383} &= -30 - 11\,x_{52} - x_{52}^{2} - 4\,x_{52} x_{503} \\
&\quad{}- 22\,x_{503} - 4\,x_{503}^{2} + x_{504} \\
r_{384} &= -306 - 35\,x_{51} - x_{51}^{2} - 16\,x_{51} x_{503} \\
&\quad{}- 4\,x_{51} x_{504} - 280\,x_{503} - 64\,x_{503}^{2} - 32\,x_{503} x_{504} \\
&\quad{}- 70\,x_{504} - 4\,x_{504}^{2} + x_{505} \\
r_{385} &= -148 - 64\,x_{51} - 32\,x_{477} - 8\,x_{482} \\
&\quad{}- 8\,x_{482}^{2} - 2\,x_{499} - 2\,x_{499}^{2} - 32\,x_{503} \\
&\quad{}- 8\,x_{504} - 2\,x_{505} + x_{506} \\
r_{386} &= -264 - 128\,x_{503} - 32\,x_{504} - 8\,x_{505} \\
&\quad{}- 2\,x_{506} - 2\,x_{506}^{2} + x_{507} \\
r_{387} &= -2 - 3\,x_{50} - x_{50}^{2} - 2\,x_{50} x_{507} \\
&\quad{}- 3\,x_{507} - x_{507}^{2} + x_{508} \\
r_{388} &= -12 - 28\,x_{507} - 16\,x_{507}^{2} - 16\,x_{507} x_{508} \\
&\quad{}- 14\,x_{508} - 4\,x_{508}^{2} + x_{509} \\
r_{389} &= -30 - 11\,x_{52} - x_{52}^{2} - 4\,x_{52} x_{509} \\
&\quad{}- 22\,x_{509} - 4\,x_{509}^{2} + x_{510} \\
r_{390} &= -306 - 35\,x_{50} - x_{50}^{2} - 16\,x_{50} x_{509} \\
&\quad{}- 4\,x_{50} x_{510} - 280\,x_{509} - 64\,x_{509}^{2} - 32\,x_{509} x_{510} \\
&\quad{}- 70\,x_{510} - 4\,x_{510}^{2} + x_{511} \\
r_{391} &= -4556 - 4320\,x_{509} - 1024\,x_{509}^{2} - 512\,x_{509} x_{510} \\
&\quad{}- 128\,x_{509} x_{511} - 1080\,x_{510} - 64\,x_{510}^{2} - 32\,x_{510} x_{511} \\
&\quad{}- 270\,x_{511} - 4\,x_{511}^{2} + x_{512} \\
r_{392} &= -90 - 76\,x_{52} - 16\,x_{52}^{2} - 16\,x_{52} x_{477} \\
&\quad{}- 16\,x_{52} x_{512} - 38\,x_{477} - 4\,x_{477}^{2} - 8\,x_{477} x_{512} \\
&\quad{}- 38\,x_{512} - 4\,x_{512}^{2} + x_{513} \\
r_{393} &= -4556 - 4320\,x_{503} - 1024\,x_{503}^{2} - 512\,x_{503} x_{504} \\
&\quad{}- 128\,x_{503} x_{505} - 1080\,x_{504} - 64\,x_{504}^{2} - 32\,x_{504} x_{505} \\
&\quad{}- 270\,x_{505} - 4\,x_{505}^{2} + x_{514} \\
r_{394} &= -682 - 70\,x_{51} - 2\,x_{51}^{2} - 32\,x_{51} x_{512} \\
&\quad{}- 8\,x_{51} x_{513} - 592\,x_{512} - 128\,x_{512}^{2} - 64\,x_{512} x_{513} \\
&\quad{}- 148\,x_{513} - 8\,x_{513}^{2} - 2\,x_{514} + x_{515} \\
r_{395} &= -x_{515} - x_{515}^{2} + x_{516} \\
r_{396} &= -42 - 52\,x_{51} - 16\,x_{51}^{2} - 8\,x_{51} x_{52} \\
&\quad{}- 16\,x_{51} x_{473} - 13\,x_{52} - x_{52}^{2} - 4\,x_{52} x_{473} \\
&\quad{}- 26\,x_{473} - 4\,x_{473}^{2} + x_{517} \\
r_{397} &= -2278 - 2264\,x_{51} - 800\,x_{51}^{2} - 800\,x_{51} x_{473} \\
&\quad{}- 160\,x_{51} x_{517} - 70\,x_{52} - 2\,x_{52}^{2} - 32\,x_{52} x_{514} \\
&\quad{}- 8\,x_{52} x_{516} - 1132\,x_{473} - 200\,x_{473}^{2} - 80\,x_{473} x_{517} \\
&\quad{}- 592\,x_{514} - 128\,x_{514}^{2} - 64\,x_{514} x_{516} - 148\,x_{516} \\
&\quad{}- 8\,x_{516}^{2} - 228\,x_{517} - 8\,x_{517}^{2} + x_{518} \\
r_{398} &= -2712 - 280\,x_{52} - 8\,x_{52}^{2} - 128\,x_{52} x_{514} \\
&\quad{}- 32\,x_{52} x_{516} - 2368\,x_{514} - 512\,x_{514}^{2} - 256\,x_{514} x_{516} \\
&\quad{}- 592\,x_{516} - 32\,x_{516}^{2} - 2\,x_{518} - 2\,x_{518}^{2} \\
&\quad{}+ x_{519} \\
r_{399} &= -2 - 3\,x_{519} - x_{519}^{2} + x_{520} \\
r_{400} &= -1 - 2\,x_{47} - x_{47}^{2} + x_{521} \\
r_{401} &= -1 - 2\,x_{48} - x_{48}^{2} + x_{522} \\
r_{402} &= -1 - 2\,x_{49} - x_{49}^{2} + x_{523} \\
r_{403} &= -16 - 16\,x_{476} - 4\,x_{476}^{2} + x_{524} \\
r_{404} &= -8 - 8\,x_{49} - 4\,x_{49} x_{476} - 4\,x_{476} \\
&\quad{}+ x_{525} \\
r_{405} &= -x_{481}^{2} + x_{526} \\
r_{406} &= -2\,x_{49} x_{481} - 2\,x_{481} + x_{527} \\
r_{407} &= -x_{484}^{2} + x_{528} \\
r_{408} &= -2\,x_{49} x_{484} - 2\,x_{484} + x_{529} \\
r_{409} &= -x_{486}^{2} + x_{530} \\
r_{410} &= -2\,x_{49} x_{486} - 2\,x_{486} + x_{531} \\
r_{411} &= -x_{488}^{2} + x_{532} \\
r_{412} &= -2\,x_{49} x_{488} - 2\,x_{488} + x_{533} \\
r_{413} &= -16 - 32\,x_{51} - 16\,x_{51}^{2} - 16\,x_{51} x_{492} \\
&\quad{}- 16\,x_{492} - 4\,x_{492}^{2} + x_{534} \\
r_{414} &= -8 - 8\,x_{49} - 8\,x_{49} x_{51} - 4\,x_{49} x_{492} \\
&\quad{}- 8\,x_{51} - 4\,x_{492} + x_{535} \\
r_{415} &= -16\,x_{494}^{2} - 16\,x_{494} x_{495} - 4\,x_{495}^{2} + x_{536} \\
r_{416} &= -8\,x_{49} x_{494} - 4\,x_{49} x_{495} - 8\,x_{494} - 4\,x_{495} \\
&\quad{}+ x_{537} \\
r_{417} &= -x_{498}^{2} + x_{538} \\
r_{418} &= -2\,x_{49} x_{498} - 2\,x_{498} + x_{539} \\
r_{419} &= -8 - 8\,x_{49} - 4\,x_{49} x_{520} - 4\,x_{520} \\
&\quad{}+ x_{540} \\
r_{420} &= -1 - 2\,x_{52} - x_{52}^{2} + x_{541} \\
r_{421} &= -1 - x_{50} - x_{50} x_{51} - x_{51} \\
&\quad{}+ x_{542} \\
r_{422} &= -2 - x_{50} - x_{50} x_{52} - x_{51} \\
&\quad{}- x_{51} x_{52} - 2\,x_{52} + x_{543} \\
r_{423} &= -x_{541}^{2} - 2\,x_{541} x_{542} - x_{542}^{2} + x_{544} \\
r_{424} &= -x_{543}^{2} + x_{545} \\
r_{425} &= -2\,x_{541} x_{543} - 2\,x_{542} x_{543} + x_{546} \\
r_{426} &= -4 - 4\,x_{53} - x_{53}^{2} + x_{547} \\
r_{427} &= -2\,x_{544} x_{546} - 2\,x_{544} x_{547} - 2\,x_{545} x_{546} - 2\,x_{545} x_{547} \\
&\quad{}+ x_{548} \\
r_{428} &= -16 - 16\,x_{520} - 4\,x_{520}^{2} - x_{521} \\
&\quad{}- x_{522} - x_{523} - x_{544}^{2} - 2\,x_{544} x_{545} \\
&\quad{}- x_{545}^{2} - x_{546}^{2} - 2\,x_{546} x_{547} - x_{547}^{2} \\
&\quad{}+ x_{549} \\
r_{429} &= -6 - 5\,x_{47} - x_{47}^{2} - 2\,x_{47} x_{49} \\
&\quad{}- 5\,x_{49} - x_{49}^{2} + x_{550} \\
r_{430} &= -16 - 32\,x_{49} - 16\,x_{49}^{2} - 16\,x_{49} x_{550} \\
&\quad{}- 16\,x_{550} - 4\,x_{550}^{2} + x_{551} \\
r_{431} &= -8 - 8\,x_{48} - 8\,x_{48} x_{49} - 4\,x_{48} x_{550} \\
&\quad{}- 8\,x_{49} - 4\,x_{550} + x_{552} \\
r_{432} &= -6 - 5\,x_{47} - x_{47}^{2} - 2\,x_{47} x_{50} \\
&\quad{}- 5\,x_{50} - x_{50}^{2} + x_{553} \\
r_{433} &= -36 - 48\,x_{47} - 16\,x_{47}^{2} - 16\,x_{47} x_{553} \\
&\quad{}- 24\,x_{553} - 4\,x_{553}^{2} + x_{554} \\
r_{434} &= -12 - 8\,x_{47} - 8\,x_{47} x_{49} - 12\,x_{49} \\
&\quad{}- 4\,x_{49} x_{553} - 4\,x_{553} + x_{555} \\
r_{435} &= -x_{521}^{2} - 2\,x_{521} x_{522} - 2\,x_{521} x_{523} - x_{521} x_{524} \\
&\quad{}- x_{522}^{2} - 2\,x_{522} x_{523} - x_{522} x_{524} - x_{523}^{2} \\
&\quad{}- x_{523} x_{524} + x_{556} \\
r_{436} &= -x_{521} x_{525} - x_{522} x_{525} - x_{523} x_{525} + x_{557} \\
r_{437} &= -x_{521} x_{556} - x_{522} x_{556} - x_{523} x_{556} - x_{526} x_{556} \\
&\quad{}+ x_{558} \\
r_{438} &= -x_{527} x_{557} + x_{559} \\
r_{439} &= -x_{527} x_{556} + x_{560} \\
r_{440} &= -x_{521} x_{557} - x_{522} x_{557} - x_{523} x_{557} - x_{526} x_{557} \\
&\quad{}+ x_{561} \\
r_{441} &= -x_{521} x_{558} - x_{521} x_{559} - x_{522} x_{558} - x_{522} x_{559} \\
&\quad{}- x_{523} x_{558} - x_{523} x_{559} - x_{528} x_{558} - x_{528} x_{559} \\
&\quad{}+ x_{562} \\
r_{442} &= -x_{529} x_{560} - x_{529} x_{561} + x_{563} \\
r_{443} &= -x_{529} x_{558} - x_{529} x_{559} + x_{564} \\
r_{444} &= -x_{521} x_{560} - x_{521} x_{561} - x_{522} x_{560} - x_{522} x_{561} \\
&\quad{}- x_{523} x_{560} - x_{523} x_{561} - x_{528} x_{560} - x_{528} x_{561} \\
&\quad{}+ x_{565} \\
r_{445} &= -x_{521} x_{562} - x_{521} x_{563} - x_{522} x_{562} - x_{522} x_{563} \\
&\quad{}- x_{523} x_{562} - x_{523} x_{563} - x_{530} x_{562} - x_{530} x_{563} \\
&\quad{}+ x_{566} \\
r_{446} &= -x_{531} x_{564} - x_{531} x_{565} + x_{567} \\
r_{447} &= -x_{531} x_{562} - x_{531} x_{563} + x_{568} \\
r_{448} &= -x_{521} x_{564} - x_{521} x_{565} - x_{522} x_{564} - x_{522} x_{565} \\
&\quad{}- x_{523} x_{564} - x_{523} x_{565} - x_{530} x_{564} - x_{530} x_{565} \\
&\quad{}+ x_{569} \\
r_{449} &= -x_{521} x_{566} - x_{521} x_{567} - x_{522} x_{566} - x_{522} x_{567} \\
&\quad{}- x_{523} x_{566} - x_{523} x_{567} - x_{532} x_{566} - x_{532} x_{567} \\
&\quad{}+ x_{570} \\
r_{450} &= -x_{533} x_{568} - x_{533} x_{569} + x_{571} \\
r_{451} &= -x_{533} x_{566} - x_{533} x_{567} + x_{572} \\
r_{452} &= -x_{521} x_{568} - x_{521} x_{569} - x_{522} x_{568} - x_{522} x_{569} \\
&\quad{}- x_{523} x_{568} - x_{523} x_{569} - x_{532} x_{568} - x_{532} x_{569} \\
&\quad{}+ x_{573} \\
r_{453} &= -x_{521} x_{570} - x_{521} x_{571} - x_{522} x_{570} - x_{522} x_{571} \\
&\quad{}- x_{523} x_{570} - x_{523} x_{571} - x_{534} x_{570} - x_{534} x_{571} \\
&\quad{}+ x_{574} \\
r_{454} &= -x_{535} x_{572} - x_{535} x_{573} + x_{575} \\
r_{455} &= -x_{535} x_{570} - x_{535} x_{571} + x_{576} \\
r_{456} &= -x_{521} x_{572} - x_{521} x_{573} - x_{522} x_{572} - x_{522} x_{573} \\
&\quad{}- x_{523} x_{572} - x_{523} x_{573} - x_{534} x_{572} - x_{534} x_{573} \\
&\quad{}+ x_{577} \\
r_{457} &= -x_{521} x_{574} - x_{521} x_{575} - x_{522} x_{574} - x_{522} x_{575} \\
&\quad{}- x_{523} x_{574} - x_{523} x_{575} - x_{536} x_{574} - x_{536} x_{575} \\
&\quad{}+ x_{578} \\
r_{458} &= -x_{537} x_{576} - x_{537} x_{577} + x_{579} \\
r_{459} &= -x_{537} x_{574} - x_{537} x_{575} + x_{580} \\
r_{460} &= -x_{521} x_{576} - x_{521} x_{577} - x_{522} x_{576} - x_{522} x_{577} \\
&\quad{}- x_{523} x_{576} - x_{523} x_{577} - x_{536} x_{576} - x_{536} x_{577} \\
&\quad{}+ x_{581} \\
r_{461} &= -x_{521} x_{578} - x_{521} x_{579} - x_{522} x_{578} - x_{522} x_{579} \\
&\quad{}- x_{523} x_{578} - x_{523} x_{579} - x_{538} x_{578} - x_{538} x_{579} \\
&\quad{}+ x_{582} \\
r_{462} &= -x_{539} x_{580} - x_{539} x_{581} + x_{583} \\
r_{463} &= -x_{539} x_{578} - x_{539} x_{579} + x_{584} \\
r_{464} &= -x_{521} x_{580} - x_{521} x_{581} - x_{522} x_{580} - x_{522} x_{581} \\
&\quad{}- x_{523} x_{580} - x_{523} x_{581} - x_{538} x_{580} - x_{538} x_{581} \\
&\quad{}+ x_{585} \\
r_{465} &= -x_{549} x_{582} - x_{549} x_{583} + x_{586} \\
r_{466} &= -x_{540} x_{584} - x_{540} x_{585} - x_{548} x_{584} - x_{548} x_{585} \\
&\quad{}+ x_{587} \\
r_{467} &= -x_{540} x_{582} - x_{540} x_{583} - x_{548} x_{582} - x_{548} x_{583} \\
&\quad{}+ x_{588} \\
r_{468} &= -x_{549} x_{584} - x_{549} x_{585} + x_{589} \\
r_{469} &= -x_{522} x_{586} - x_{522} x_{587} - x_{551} x_{586} - x_{551} x_{587} \\
&\quad{}+ x_{590} \\
r_{470} &= -x_{552} x_{588} - x_{552} x_{589} + x_{591} \\
r_{471} &= -x_{552} x_{586} - x_{552} x_{587} + x_{592} \\
r_{472} &= -x_{522} x_{588} - x_{522} x_{589} - x_{551} x_{588} - x_{551} x_{589} \\
&\quad{}+ x_{593} \\
r_{473} &= -1 - x_{4} - x_{4} x_{49} - x_{49} \\
&\quad{}+ x_{594} \\
r_{474} &= -1 - x_{55} - x_{55} x_{594} - x_{594} \\
&\quad{}+ x_{595} \\
r_{475} &= -1 - x_{4} - x_{4} x_{47} - x_{47} \\
&\quad{}+ x_{596} \\
r_{476} &= -1 - x_{4} - x_{4} x_{48} - x_{48} \\
&\quad{}+ x_{597} \\
r_{477} &= -1 - x_{54} - x_{54} x_{596} - x_{596} \\
&\quad{}+ x_{598} \\
r_{478} &= -x_{597} x_{598} - x_{598} + x_{599} \\
r_{479} &= -x_{2} - x_{2} x_{58} + x_{600} \\
r_{480} &= -1 - x_{56} - x_{56} x_{600} - x_{600} \\
&\quad{}+ x_{601} \\
r_{481} &= -1 - x_{2} - x_{2} x_{57} - x_{57} \\
&\quad{}- x_{57} x_{600} - x_{600} + x_{602} \\
r_{482} &= 2 - 4\,x_{1} + 2\,x_{1}^{2} - 2\,x_{1} x_{54} \\
&\quad{}- 2\,x_{1} x_{55} - 2\,x_{1} x_{601} - 2\,x_{1} x_{602} + 2\,x_{54} \\
&\quad{}+ x_{54}^{2} + 2\,x_{54} x_{601} + 2\,x_{55} + x_{55}^{2} \\
&\quad{}+ 2\,x_{55} x_{602} + x_{195} x_{603} - x_{195} x_{604} + x_{522} x_{590} \\
&\quad{}+ x_{522} x_{591} - x_{522} x_{592} - x_{522} x_{593} + x_{523} x_{590} \\
&\quad{}+ x_{523} x_{591} - x_{523} x_{592} - x_{523} x_{593} + x_{554} x_{590} \\
&\quad{}+ x_{554} x_{591} - x_{554} x_{592} - x_{554} x_{593} - x_{555} x_{590} \\
&\quad{}- x_{555} x_{591} + x_{555} x_{592} + x_{555} x_{593} + x_{595}^{2} \\
&\quad{}- 2\,x_{595} x_{599} + x_{599}^{2} + 2\,x_{601} + x_{601}^{2} \\
&\quad{}+ 2\,x_{602} + x_{602}^{2} \\
r_{483} &= -15 - 2\,x_{0} + x_{5} - x_{47} \\
&\quad{}- x_{48} - x_{49} - x_{50} - x_{51} \\
&\quad{}- x_{52} - x_{53} - x_{54} - x_{55} \\
&\quad{}- x_{56} - x_{57} - x_{92} - x_{95} \\
&\quad{}- x_{136} - x_{144} \\
r_{484} &= -69 - 7\,x_{0} - 2\,x_{2} - 3\,x_{3} \\
&\quad{}- x_{5} - 19\,x_{6} + x_{7} - x_{46} \\
&\quad{}- x_{47} - x_{48} - x_{49} - x_{50} \\
&\quad{}- x_{51} - x_{52} - x_{53} - x_{54} \\
&\quad{}- x_{55} - x_{56} - x_{57} - 3\,x_{92} \\
&\quad{}- x_{93} - 2\,x_{95} - x_{135} - x_{136} \\
&\quad{}- 2\,x_{144} \\
r_{485} &= -x_{8} x_{11} + x_{605} \\
r_{486} &= -1 + x_{10}^{2} + x_{11}^{2} - x_{605}^{2} \\
r_{487} &= -x_{7} x_{9} + x_{606} \\
r_{488} &= -x_{7} x_{606} - x_{606} + x_{607} \\
r_{489} &= -1 + x_{8}^{2} + x_{606}^{2} - x_{607}^{2} \\
r_{490} &= -x_{8} - x_{10} x_{13} + x_{12}
\end{align*}

\endgroup

\clearpage
\section{Integer input-guard index set}\label{app:integer}

Keep the 608 coordinates and add three fresh integer coordinates for
each following index, in ascending order. Together with
\eqref{eq:integer}, this defines the entire integer polynomial.
\begin{align*}
I=\{&0, 1, 2, 3, 4, 5, 6, 7, 8, 9, 10, 11, 12, 13, 14, 15, 16, 17, 20, 21,\\
&23, 24, 25, 26, 27, 28, 29, 30, 31, 34, 35, 37, 38, 39, 40, 41, 42, 43, 44, 46,\\
&47, 48, 49, 50, 51, 52, 53, 54, 55, 56, 57, 58, 93, 96, 97, 98, 99, 100, 101, 104,\\
&105, 107, 108, 109, 110, 111, 112, 113, 114, 115, 118, 119, 121, 122, 123, 124, 125, 126, 127, 128,\\
&130, 131, 132, 133, 134, 136, 137, 138, 139, 140, 141, 142, 145, 146, 147, 148, 149, 150, 153, 154,\\
&156, 157, 158, 159, 160, 161, 162, 163, 164, 165, 166, 168, 169, 170, 171, 172, 173, 174, 175, 176,\\
&177, 178, 181, 182, 184, 185, 186, 187, 188, 189, 190, 191, 194, 195, 196, 197, 198, 199, 200, 201,\\
&202, 203, 206, 207, 209, 210, 211, 212, 213, 214, 215, 216, 217, 220, 221, 223, 224, 225, 226, 227,\\
&228, 229, 230, 232, 233, 234, 235, 236, 237, 238, 241, 242, 244, 245, 246, 247, 248, 249, 250, 251,\\
&253, 254, 255, 256, 257, 258, 259, 262, 263, 265, 266, 267, 268, 269, 270, 271, 272, 274, 275, 276,\\
&277, 278, 279, 280, 283, 284, 286, 287, 288, 289, 290, 291, 292, 293, 294, 298, 299, 301, 302, 303,\\
&304, 305, 306, 307, 308, 310, 311, 312, 313, 314, 315, 316, 319, 320, 322, 323, 324, 325, 326, 327,\\
&328, 329, 331, 332, 333, 334, 335, 336, 337, 340, 341, 343, 344, 345, 346, 347, 348, 349, 350, 352,\\
&353, 354, 355, 356, 357, 358, 361, 362, 364, 365, 366, 367, 368, 369, 370, 371, 373, 374, 375, 376,\\
&377, 378, 379, 382, 383, 385, 386, 387, 388, 389, 390, 391, 392, 393, 396, 397, 399, 400, 401, 402,\\
&403, 404, 405, 406, 408, 409, 410, 411, 412, 413, 414, 417, 418, 420, 421, 422, 423, 424, 425, 426,\\
&427, 429, 430, 431, 432, 433, 434, 435, 438, 439, 441, 442, 443, 444, 445, 446, 447, 448, 450, 451,\\
&452, 453, 454, 455, 456, 459, 460, 462, 463, 464, 465, 466, 467, 468, 469, 471, 472, 603, 604\}.
\end{align*}


\end{document}
